% figs.tex —— 由 ch9/make_figs.py 產生，請勿手動修改。
% 重新產生：python ch9/make_figs.py（第 9 章，動畫 W901–W906）

\newcommand{\NineFigMs}{%
\begin{tikzpicture}[x=4.2cm,y=2.6cm,>=Stealth,font=\scriptsize,baseline=(current bounding box.center)]
\draw[->] (0,0) -- (1.3,0) node[right] {$T/T_c$};
\draw[->] (0,0) -- (0,1.14) node[above] {$M_s$};
\node[below] at (0,-0.03) {$0$};
\foreach \x in {0.5,1} {\draw (\x,0) -- (\x,-0.03) node[below] {$\x$};}
\foreach \y in {0.5,1} {\draw (0,\y) -- (-0.02,\y) node[left] {$\y$};}
\draw[gray!50,densely dotted] (1,0) -- (1,1.05);
\draw[cGauss,thick] plot coordinates {(0.0000,1.0000) (0.0500,1.0000) (0.0947,1.0000) (0.1395,1.0000) (0.1842,1.0000) (0.2289,0.9997) (0.2737,0.9986) (0.3184,0.9962) (0.3632,0.9915) (0.4079,0.9841) (0.4526,0.9732) (0.4974,0.9585) (0.5421,0.9394) (0.5868,0.9154) (0.6316,0.8859) (0.6763,0.8503) (0.7211,0.8076) (0.7658,0.7564) (0.8105,0.6949) (0.8553,0.6198) (0.9000,0.5254) (0.9100,0.5006) (0.9439,0.4009) (0.9650,0.3193) (0.9782,0.2534) (0.9864,0.2008) (0.9915,0.1588) (0.9947,0.1255) (0.9967,0.0992) (0.9980,0.0784) (0.9987,0.0619) (0.9992,0.0489) (0.9995,0.0386) (0.9997,0.0305) (0.9998,0.0240) (0.9999,0.0190) (0.9999,0.0150) (1.0000,0.0118) (1.0000,0.0093) (1.0000,0.0074) (1.0000,0.0058) (1.0000,0.0046) (1.0000,0.0036) (1.0000,0.0029) (1.0000,0.0023) (1.0000,0.0018) (1.0000,0.0014) (1.0000,0.0011) (1.0000,0.0009) (1.0000,0.0007) (1.0000,0.0005) (1.0000,0.0000) (1.2000,0.0000)};
\fill[cGauss] (0.5,0.9575) circle (1.4pt);
\fill[cGauss] (0.9,0.5254) circle (1.4pt);
\fill[cGauss] (0.99,0.1725) circle (1.4pt);
\draw[cBar,thick] plot coordinates {(0.0000,1.0000) (0.0500,1.0000) (0.0947,1.0000) (0.1395,1.0000) (0.1842,1.0000) (0.2289,1.0000) (0.2737,1.0000) (0.3184,1.0000) (0.3632,0.9999) (0.4079,0.9996) (0.4526,0.9991) (0.4974,0.9981) (0.5421,0.9964) (0.5868,0.9938) (0.6316,0.9900) (0.6763,0.9844) (0.7211,0.9765) (0.7658,0.9656) (0.8105,0.9503) (0.8553,0.9286) (0.9000,0.8959) (0.9100,0.8863) (0.9439,0.8420) (0.9650,0.7975) (0.9782,0.7540) (0.9864,0.7120) (0.9915,0.6720) (0.9947,0.6339) (0.9967,0.5978) (0.9980,0.5636) (0.9987,0.5314) (0.9992,0.5009) (0.9995,0.4722) (0.9997,0.4451) (0.9998,0.4196) (0.9999,0.3955) (0.9999,0.3728) (1.0000,0.3514) (1.0000,0.3313) (1.0000,0.3123) (1.0000,0.2944) (1.0000,0.2775) (1.0000,0.2615) (1.0000,0.2465) (1.0000,0.2324) (1.0000,0.2190) (1.0000,0.2065) (1.0000,0.1946) (1.0000,0.1835) (1.0000,0.1729) (1.0000,0.1630) (1.0000,0.0000) (1.2000,0.0000)};
\fill[cBar] (0.5,0.9980) circle (1.4pt);
\fill[cBar] (0.9,0.8959) circle (1.4pt);
\fill[cBar] (0.99,0.6859) circle (1.4pt);
\draw[cBar,thick] (0.05,0.34) -- (0.13,0.34) node[right] {二維 Ising，$\beta=\frac18$};
\draw[cGauss,thick] (0.05,0.2) -- (0.13,0.2) node[right] {平均場，$\beta=\frac12$};
\end{tikzpicture}}

\newcommand{\NineTabFlory}{%
\begin{tabular}{@{}lcccc@{}}
\toprule
$d$ & 1 & 2 & 3 & 4\\
\midrule
Flory $p=\frac6{d+2}$ & $2$ & $\frac{3}{2}$ & $\frac{6}{5}$ & $1$\\
已知的 $2\nu$ & $2$ & $\frac{3}{2}$ & 1.176 & $1$\\
\bottomrule
\end{tabular}}

\newcommand{\NineTabMs}{%
\begin{tabular}{@{}lcccc@{}}
\toprule
$T/T_c$ & 0.5 & 0.9 & 0.99 & 斜率\\
\midrule
平均場 & 0.9575 & 0.5254 & 0.1725 & 0.5000\\
Onsager & 0.9980 & 0.8959 & 0.6859 & 0.1250\\
\bottomrule
\end{tabular}}

\newcommand{\NineFigPM}{%
\begin{tikzpicture}[x=0.6cm,y=4.0cm,font=\scriptsize,baseline=(current bounding box.center)]
\draw (-5,0) -- (5,0);
\node[anchor=south west] at (-5,0.5) {$P(M)$};
\fill[cBar!80] (-4.56,0) rectangle (-4.04,0.4504);
\fill[cBar!80] (-2.56,0) rectangle (-2.04,0.0330);
\fill[cBar!80] (-0.56,0) rectangle (-0.04,0.0333);
\fill[cBar!80] (1.44,0) rectangle (1.96,0.0330);
\fill[cBar!80] (3.44,0) rectangle (3.96,0.4504);
\fill[cWalk!80] (-3.96,0) rectangle (-3.44,0.1283);
\fill[cWalk!80] (-1.96,0) rectangle (-1.44,0.2306);
\fill[cWalk!80] (0.04,0) rectangle (0.56,0.2823);
\fill[cWalk!80] (2.04,0) rectangle (2.56,0.2306);
\fill[cWalk!80] (4.04,0) rectangle (4.56,0.1283);
\node[cBar,above] at (3.7,0.4504) {0.4504};
\node[cWalk,above] at (0.3,0.2823) {0.2823};
\foreach \m in {-4,-2,0,2,4} {\node[below] at (\m,0) {$\m$};}
\node[below] at (5,0) {$M$};
\fill[cBar!80] (-2.6,0.47) rectangle (-2.3,0.5); \node[anchor=west] at (-2.25,0.485) {$K=1$（$T=1$）};
\fill[cWalk!80] (-2.6,0.4) rectangle (-2.3,0.43); \node[anchor=west] at (-2.25,0.415) {$K=0.2$（$T=5$）};
\end{tikzpicture}}

\newcommand{\NineFigTrace}{%
\begin{tikzpicture}[x=1.2cm,y=0.42cm,font=\scriptsize,baseline=(current bounding box.center)]
\draw[gray!50] (0,-1.0) rectangle (4,1.0);
\draw[gray!50,densely dotted] (0,-0.0) -- (4,-0.0);
\draw[cBar,thin] plot coordinates {(0.001,0.812) (0.021,-0.031) (0.041,0.156) (0.061,0.469) (0.081,-0.219) (0.101,0.094) (0.121,0.094) (0.141,0.312) (0.161,-0.250) (0.181,0.250) (0.201,-0.250) (0.221,0.406) (0.242,-0.094) (0.262,0.125) (0.282,-0.500) (0.302,0.281) (0.322,0.344) (0.342,0.375) (0.362,-0.344) (0.382,-0.562) (0.402,-0.094) (0.422,-0.625) (0.442,0.250) (0.462,-0.500) (0.482,-0.719) (0.502,-0.812) (0.522,-0.375) (0.542,0.406) (0.562,0.375) (0.582,0.219) (0.602,-0.031) (0.622,-0.562) (0.642,-0.094) (0.662,-0.188) (0.683,0.469) (0.703,-0.281) (0.723,0.406) (0.743,-0.094) (0.763,0.375) (0.783,0.094) (0.803,0.531) (0.823,0.656) (0.843,0.750) (0.863,0.000) (0.883,-0.188) (0.903,-0.281) (0.923,0.094) (0.943,0.469) (0.963,-0.094) (0.983,-0.562) (1.003,-0.594) (1.023,-0.031) (1.043,0.000) (1.063,-0.062) (1.083,0.250) (1.103,-0.062) (1.124,-0.812) (1.144,0.594) (1.164,-0.688) (1.184,-0.219) (1.204,-0.125) (1.224,-0.281) (1.244,0.688) (1.264,0.656) (1.284,-0.656) (1.304,-0.219) (1.324,-0.125) (1.344,0.062) (1.364,-0.438) (1.384,0.094) (1.404,0.000) (1.424,-0.094) (1.444,-0.656) (1.464,0.375) (1.484,-0.469) (1.504,-0.031) (1.524,-0.688) (1.544,0.031) (1.565,0.094) (1.585,0.500) (1.605,-0.312) (1.625,0.250) (1.645,-0.125) (1.665,0.500) (1.685,-0.500) (1.705,0.094) (1.725,-0.750) (1.745,0.750) (1.765,0.438) (1.785,-0.125) (1.805,0.562) (1.825,-0.531) (1.845,-0.500) (1.865,-0.062) (1.885,0.312) (1.905,0.625) (1.925,0.156) (1.945,0.312) (1.965,0.438) (1.985,-0.156) (2.006,-0.094) (2.026,-0.312) (2.046,0.250) (2.066,0.469) (2.086,-0.312) (2.106,0.844) (2.126,0.156) (2.146,0.219) (2.166,0.125) (2.186,0.688) (2.206,-0.031) (2.226,0.000) (2.246,0.031) (2.266,0.344) (2.286,-0.125) (2.306,-0.312) (2.326,-0.219) (2.346,-0.344) (2.366,0.906) (2.386,-0.062) (2.406,-0.125) (2.426,0.406) (2.447,0.406) (2.467,0.688) (2.487,0.281) (2.507,-0.125) (2.527,0.500) (2.547,-0.469) (2.567,-0.188) (2.587,0.000) (2.607,-0.125) (2.627,-0.312) (2.647,0.250) (2.667,0.562) (2.687,-0.750) (2.707,0.469) (2.727,0.531) (2.747,-0.719) (2.767,-0.312) (2.787,-0.156) (2.807,-0.406) (2.827,-0.406) (2.847,-0.781) (2.867,-0.062) (2.887,0.188) (2.908,-0.562) (2.928,-0.156) (2.948,-0.125) (2.968,-0.125) (2.988,0.344) (3.008,0.656) (3.028,0.406) (3.048,0.562) (3.068,0.406) (3.088,0.812) (3.108,0.281) (3.128,-0.375) (3.148,-0.156) (3.168,0.438) (3.188,-0.344) (3.208,-0.344) (3.228,-0.500) (3.248,0.156) (3.268,0.125) (3.288,-0.281) (3.308,0.656) (3.328,0.625) (3.349,0.719) (3.369,-0.469) (3.389,0.625) (3.409,-0.156) (3.429,0.000) (3.449,-0.281) (3.469,0.344) (3.489,-0.562) (3.509,-0.375) (3.529,-0.781) (3.549,-0.938) (3.569,-0.656) (3.589,-0.375) (3.609,0.469) (3.629,0.000) (3.649,0.000) (3.669,0.312) (3.689,-0.406) (3.709,0.031) (3.729,-0.531) (3.749,0.688) (3.769,-0.031) (3.790,0.844) (3.810,-0.406) (3.830,0.031) (3.850,0.031) (3.870,0.156) (3.890,0.062) (3.910,-0.438) (3.930,-0.531) (3.950,-0.594) (3.970,-0.188) (3.990,0.344)};
\node[anchor=west,align=left,inner sep=2pt] at (4.05,-0.0) {$L=8$，$T=3.0$\\（高溫）};
\draw[gray!50] (0,-3.5) rectangle (4,-1.5);
\draw[gray!50,densely dotted] (0,-2.5) -- (4,-2.5);
\draw[cBar,thin] plot coordinates {(0.001,-1.500) (0.021,-1.875) (0.041,-1.500) (0.061,-3.500) (0.081,-3.500) (0.101,-3.375) (0.121,-3.500) (0.141,-3.375) (0.161,-3.500) (0.181,-3.500) (0.201,-3.500) (0.221,-3.500) (0.242,-3.375) (0.262,-3.500) (0.282,-3.500) (0.302,-3.500) (0.322,-3.500) (0.342,-3.125) (0.362,-3.375) (0.382,-3.500) (0.402,-1.500) (0.422,-1.500) (0.442,-1.500) (0.462,-1.500) (0.482,-1.500) (0.502,-1.500) (0.522,-3.500) (0.542,-1.500) (0.562,-1.500) (0.582,-1.500) (0.602,-1.500) (0.622,-2.125) (0.642,-1.500) (0.662,-1.750) (0.683,-1.750) (0.703,-1.500) (0.723,-1.500) (0.743,-1.500) (0.763,-1.500) (0.783,-3.125) (0.803,-3.500) (0.823,-3.500) (0.843,-3.375) (0.863,-3.375) (0.883,-3.250) (0.903,-3.375) (0.923,-3.500) (0.943,-1.750) (0.963,-1.500) (0.983,-1.500) (1.003,-3.375) (1.023,-3.375) (1.043,-3.500) (1.063,-3.500) (1.083,-3.250) (1.103,-3.500) (1.124,-3.375) (1.144,-3.250) (1.164,-3.500) (1.184,-1.500) (1.204,-1.875) (1.224,-1.500) (1.244,-1.625) (1.264,-1.875) (1.284,-3.375) (1.304,-3.500) (1.324,-3.500) (1.344,-3.500) (1.364,-3.375) (1.384,-1.500) (1.404,-1.625) (1.424,-1.500) (1.444,-1.500) (1.464,-1.625) (1.484,-1.500) (1.504,-1.500) (1.524,-1.500) (1.544,-1.625) (1.565,-1.625) (1.585,-1.500) (1.605,-1.500) (1.625,-1.500) (1.645,-1.500) (1.665,-1.500) (1.685,-1.500) (1.705,-1.625) (1.725,-1.625) (1.745,-1.625) (1.765,-1.500) (1.785,-1.500) (1.805,-1.500) (1.825,-1.500) (1.845,-3.500) (1.865,-3.500) (1.885,-3.375) (1.905,-3.500) (1.925,-3.125) (1.945,-3.375) (1.965,-3.500) (1.985,-3.500) (2.006,-3.500) (2.026,-3.500) (2.046,-3.500) (2.066,-3.500) (2.086,-3.375) (2.106,-3.125) (2.126,-3.375) (2.146,-3.250) (2.166,-3.500) (2.186,-3.500) (2.206,-3.500) (2.226,-3.500) (2.246,-3.125) (2.266,-3.375) (2.286,-3.500) (2.306,-3.000) (2.326,-2.750) (2.346,-3.375) (2.366,-3.500) (2.386,-3.500) (2.406,-3.500) (2.426,-3.500) (2.447,-3.500) (2.467,-3.500) (2.487,-3.375) (2.507,-3.500) (2.527,-3.500) (2.547,-3.500) (2.567,-3.500) (2.587,-3.375) (2.607,-3.500) (2.627,-3.500) (2.647,-3.375) (2.667,-3.500) (2.687,-3.250) (2.707,-3.500) (2.727,-3.375) (2.747,-1.500) (2.767,-3.125) (2.787,-3.500) (2.807,-3.250) (2.827,-3.500) (2.847,-3.500) (2.867,-3.500) (2.887,-3.500) (2.908,-3.500) (2.928,-3.500) (2.948,-3.375) (2.968,-2.125) (2.988,-3.500) (3.008,-3.375) (3.028,-3.500) (3.048,-3.250) (3.068,-1.500) (3.088,-1.500) (3.108,-1.500) (3.128,-1.500) (3.148,-1.500) (3.168,-3.000) (3.188,-3.500) (3.208,-1.625) (3.228,-1.625) (3.248,-1.500) (3.268,-1.500) (3.288,-1.500) (3.308,-1.500) (3.328,-1.500) (3.349,-2.000) (3.369,-1.625) (3.389,-1.500) (3.409,-1.500) (3.429,-1.625) (3.449,-1.500) (3.469,-1.500) (3.489,-1.500) (3.509,-2.250) (3.529,-3.500) (3.549,-3.500) (3.569,-3.500) (3.589,-3.375) (3.609,-1.500) (3.629,-2.000) (3.649,-1.500) (3.669,-1.500) (3.689,-1.500) (3.709,-1.875) (3.729,-3.375) (3.749,-1.500) (3.769,-1.500) (3.790,-1.500) (3.810,-1.500) (3.830,-1.625) (3.850,-1.625) (3.870,-3.500) (3.890,-3.500) (3.910,-3.500) (3.930,-1.625) (3.950,-1.500) (3.970,-1.500) (3.990,-1.500)};
\node[anchor=west,align=left,inner sep=2pt] at (4.05,-2.5) {$L=4$，$T=2.0$\\翻 30 次};
\draw[gray!50] (0,-6.0) rectangle (4,-4.0);
\draw[gray!50,densely dotted] (0,-5.0) -- (4,-5.0);
\draw[cBar,thin] plot coordinates {(0.001,-4.000) (0.021,-4.000) (0.041,-4.000) (0.061,-4.062) (0.081,-4.031) (0.101,-4.000) (0.121,-4.031) (0.141,-4.000) (0.161,-4.188) (0.181,-4.031) (0.201,-4.062) (0.221,-4.031) (0.242,-4.125) (0.262,-4.000) (0.282,-4.031) (0.302,-4.094) (0.322,-4.000) (0.342,-4.156) (0.362,-4.031) (0.382,-4.156) (0.402,-4.031) (0.422,-4.062) (0.442,-4.000) (0.462,-4.188) (0.482,-4.031) (0.502,-4.094) (0.522,-5.875) (0.542,-5.938) (0.562,-5.938) (0.582,-5.938) (0.602,-6.000) (0.622,-6.000) (0.642,-5.906) (0.662,-5.781) (0.683,-6.000) (0.703,-5.812) (0.723,-5.938) (0.743,-5.969) (0.763,-5.938) (0.783,-5.969) (0.803,-5.875) (0.823,-5.844) (0.843,-5.969) (0.863,-5.875) (0.883,-5.938) (0.903,-5.844) (0.923,-5.812) (0.943,-6.000) (0.963,-6.000) (0.983,-5.969) (1.003,-6.000) (1.023,-5.938) (1.043,-5.938) (1.063,-6.000) (1.083,-5.875) (1.103,-5.969) (1.124,-5.875) (1.144,-5.938) (1.164,-5.719) (1.184,-5.875) (1.204,-5.938) (1.224,-5.969) (1.244,-6.000) (1.264,-5.719) (1.284,-5.938) (1.304,-5.969) (1.324,-5.906) (1.344,-5.656) (1.364,-5.906) (1.384,-5.969) (1.404,-5.969) (1.424,-6.000) (1.444,-5.875) (1.464,-5.969) (1.484,-5.875) (1.504,-5.688) (1.524,-5.969) (1.544,-5.875) (1.565,-5.875) (1.585,-5.938) (1.605,-5.938) (1.625,-5.844) (1.645,-5.531) (1.665,-4.031) (1.685,-4.156) (1.705,-4.031) (1.725,-4.125) (1.745,-4.062) (1.765,-4.094) (1.785,-4.000) (1.805,-4.031) (1.825,-4.031) (1.845,-4.000) (1.865,-4.062) (1.885,-4.062) (1.905,-4.344) (1.925,-4.062) (1.945,-4.094) (1.965,-4.031) (1.985,-4.062) (2.006,-4.000) (2.026,-4.062) (2.046,-4.000) (2.066,-4.406) (2.086,-4.031) (2.106,-4.031) (2.126,-4.312) (2.146,-4.031) (2.166,-4.094) (2.186,-4.031) (2.206,-4.156) (2.226,-4.094) (2.246,-4.062) (2.266,-4.062) (2.286,-4.031) (2.306,-4.031) (2.326,-4.000) (2.346,-4.062) (2.366,-4.000) (2.386,-4.031) (2.406,-4.031) (2.426,-4.000) (2.447,-4.156) (2.467,-4.094) (2.487,-4.219) (2.507,-4.031) (2.527,-4.000) (2.547,-4.156) (2.567,-4.188) (2.587,-4.062) (2.607,-4.062) (2.627,-4.188) (2.647,-4.031) (2.667,-4.000) (2.687,-4.094) (2.707,-4.125) (2.727,-4.062) (2.747,-4.094) (2.767,-4.000) (2.787,-4.031) (2.807,-4.031) (2.827,-4.094) (2.847,-4.312) (2.867,-4.031) (2.887,-4.000) (2.908,-4.125) (2.928,-4.062) (2.948,-4.062) (2.968,-4.000) (2.988,-4.094) (3.008,-4.219) (3.028,-4.031) (3.048,-4.219) (3.068,-4.281) (3.088,-4.000) (3.108,-4.094) (3.128,-4.062) (3.148,-4.094) (3.168,-4.188) (3.188,-4.031) (3.208,-4.094) (3.228,-4.062) (3.248,-4.125) (3.268,-4.031) (3.288,-4.062) (3.308,-4.000) (3.328,-4.000) (3.349,-4.000) (3.369,-4.062) (3.389,-4.031) (3.409,-4.062) (3.429,-4.031) (3.449,-4.031) (3.469,-4.156) (3.489,-4.062) (3.509,-4.344) (3.529,-4.062) (3.549,-4.188) (3.569,-4.094) (3.589,-4.281) (3.609,-4.062) (3.629,-4.031) (3.649,-4.000) (3.669,-4.000) (3.689,-4.094) (3.709,-4.062) (3.729,-4.031) (3.749,-4.094) (3.769,-4.000) (3.790,-4.156) (3.810,-4.062) (3.830,-4.000) (3.850,-4.031) (3.870,-4.062) (3.890,-4.156) (3.910,-4.062) (3.930,-4.125) (3.950,-4.219) (3.970,-4.062) (3.990,-4.156)};
\node[anchor=west,align=left,inner sep=2pt] at (4.05,-5.0) {$L=8$，$T=2.0$\\翻 2 次};
\node[left,inner sep=1pt] at (0,1) {$1$}; \node[left,inner sep=1pt] at (0,0) {$m$}; \node[left,inner sep=1pt] at (0,-1) {$-1$};
\foreach \x/\l in {0/0,2/2000,4/4000} {\node[below] at (\x,-6.0) {\l};}
\node[below] at (3.0,-6.0) {sweep};
\end{tikzpicture}}

\newcommand{\NineTabFlip}{%
\begin{tabular}{@{}lccc@{}}
\toprule
$L$（$T=2.0$） & 4 & 6 & 8\\
\midrule
平均翻轉等待（sweep） & 133 & 608 & 2151\\
\bottomrule
\end{tabular}}

\newcommand{\NineTabField}{%
\begin{tabular}{@{}lcc@{}}
\toprule
$\langle m\rangle$（$K=1$） & $H=0.1$ & $H=0.01$\\
\midrule
4 個自旋一圈 & 0.3506 & 0.0367\\
$4\times4$ 週期 & 0.9210 & 0.1584\\
\multicolumn{3}{@{}l@{}}{無限大格子：$M_s=0.9993$}\\
\bottomrule
\end{tabular}}

\newcommand{\NineTabNineOne}{%
\begin{tabular}{@{}llccc@{}}
\toprule
量 & 行為 & 指數 & 平均場 & 三維（書上）\\
\midrule
比熱 $T\gtrless T_c$ & $|T-T_c|^{-\alpha}$ & $\alpha,\alpha'$ & 0 & {\color{cWrong}0.12}$\to${\color{cExact}$\approx$0.19}\\
磁化 $h=0$ & $|T_c-T|^{\beta}$ & $\beta$ & $\frac12$ & 0.308\\
磁化率 $T>T_c$ & $|T-T_c|^{-\gamma}$ & $\gamma$ & 1 & 1.19\\
磁化率 $T<T_c$ & $|T-T_c|^{-\gamma'}$ & $\gamma'$ & 1 & 無定義\\
磁化 $T=T_c$ & $|h|^{1/\delta}$ & $\delta$ & 3 & 4.865\\
關聯長度（{\color{cWrong}$\xi$}$\to${\color{cExact}$\xi^2$}） & $|T-T_c|^{-\nu}$ & $\nu,\nu'$ & $\frac12$ & 0.602\\
$C(T_c,0;\rho)$ & $\rho^{-d+2-\eta}$ & $\eta$ & 0 & 0.023\\
\bottomrule
\end{tabular}}

\newcommand{\NineTabLaws}{%
\begin{tabular}{@{}lcccc@{}}
\toprule
 & 平均場 & 二維 Ising & 二維 $n=0$ & 書上三維\\
\midrule
$\alpha$ & $0$ & $0$ & $\frac{1}{2}$ & 0.12\\
$\beta$ & $\frac{1}{2}$ & $\frac{1}{8}$ & $\frac{5}{64}$ & 0.308\\
$\gamma$ & $1$ & $\frac{7}{4}$ & $\frac{43}{32}$ & 1.19\\
$\delta$ & $3$ & $15$ & $\frac{91}{5}$ & 4.865\\
$\nu$ & $\frac{1}{2}$ & $1$ & $\frac{3}{4}$ & 0.602\\
$\eta$ & $0$ & $\frac{1}{4}$ & $\frac{5}{24}$ & 0.023\\
$d$ & 4 & 2 & 2 & 3\\
\midrule
$\alpha+2\beta+\gamma\overset?=2$ & $2$ {\color{cExact}$\checkmark$} & $2$ {\color{cExact}$\checkmark$} & $2$ {\color{cExact}$\checkmark$} & 1.926 {\color{cWrong}$\times$}\\
$(2-\eta)\nu\overset?=\gamma$ & $1$ {\color{cExact}$\checkmark$} & $\frac{7}{4}$ {\color{cExact}$\checkmark$} & $\frac{43}{32}$ {\color{cExact}$\checkmark$} & 1.190 {\color{cExact}$\checkmark$}\\
$2-d\nu\overset?=\alpha$ & $0$ {\color{cExact}$\checkmark$} & $0$ {\color{cExact}$\checkmark$} & $\frac{1}{2}$ {\color{cExact}$\checkmark$} & 0.194 {\color{cWrong}$\times$}\\
$\frac{d+2-\eta}{d-2+\eta}\overset?=\delta$ & $3$ {\color{cExact}$\checkmark$} & $15$ {\color{cExact}$\checkmark$} & $\frac{91}{5}$ {\color{cExact}$\checkmark$} & 4.865 {\color{cExact}$\checkmark$}\\
$\beta(\delta-1)\overset?=\gamma$ & $1$ {\color{cExact}$\checkmark$} & $\frac{7}{4}$ {\color{cExact}$\checkmark$} & $\frac{43}{32}$ {\color{cExact}$\checkmark$} & 1.190 {\color{cExact}$\checkmark$}\\
\bottomrule
\end{tabular}}

\newcommand{\NineFigChain}{%
\begin{tikzpicture}[x=0.33cm,y=1.55cm,font=\scriptsize,baseline=(current bounding box.center)]
\draw (-8.8,0) -- (8.8,0) node[right] {$r$};
\fill[cBar!70] (-8.34,0) rectangle (-7.66,0.1132);
\fill[cBar!70] (-7.34,0) rectangle (-6.66,0.1486);
\fill[cBar!70] (-6.34,0) rectangle (-5.66,0.1951);
\fill[cBar!70] (-5.34,0) rectangle (-4.66,0.2562);
\fill[cBar!70] (-4.34,0) rectangle (-3.66,0.3364);
\fill[cBar!70] (-3.34,0) rectangle (-2.66,0.4417);
\fill[cBar!70] (-2.34,0) rectangle (-1.66,0.5800);
\fill[cBar!70] (-1.34,0) rectangle (-0.66,0.7616);
\fill[cBar!70] (-0.34,0) rectangle (0.34,1.0000);
\fill[cBar!70] (0.66,0) rectangle (1.34,0.7616);
\fill[cBar!70] (1.66,0) rectangle (2.34,0.5800);
\fill[cBar!70] (2.66,0) rectangle (3.34,0.4417);
\fill[cBar!70] (3.66,0) rectangle (4.34,0.3364);
\fill[cBar!70] (4.66,0) rectangle (5.34,0.2562);
\fill[cBar!70] (5.66,0) rectangle (6.34,0.1951);
\fill[cBar!70] (6.66,0) rectangle (7.34,0.1486);
\fill[cBar!70] (7.66,0) rectangle (8.34,0.1132);
\draw[cExact,densely dashed,thick] plot coordinates {(-8.400,0.1015) (-8.100,0.1101) (-7.800,0.1195) (-7.500,0.1297) (-7.200,0.1407) (-6.900,0.1527) (-6.600,0.1657) (-6.300,0.1798) (-6.000,0.1951) (-5.700,0.2118) (-5.400,0.2298) (-5.100,0.2493) (-4.800,0.2706) (-4.500,0.2936) (-4.200,0.3186) (-3.900,0.3457) (-3.600,0.3752) (-3.300,0.4071) (-3.000,0.4417) (-2.700,0.4794) (-2.400,0.5202) (-2.100,0.5644) (-1.800,0.6125) (-1.500,0.6646) (-1.200,0.7212) (-0.900,0.7826) (-0.600,0.8492) (-0.300,0.9215) (0.000,1.0000) (0.300,0.9215) (0.600,0.8492) (0.900,0.7826) (1.200,0.7212) (1.500,0.6646) (1.800,0.6125) (2.100,0.5644) (2.400,0.5202) (2.700,0.4794) (3.000,0.4417) (3.300,0.4071) (3.600,0.3752) (3.900,0.3457) (4.200,0.3186) (4.500,0.2936) (4.800,0.2706) (5.100,0.2493) (5.400,0.2298) (5.700,0.2118) (6.000,0.1951) (6.300,0.1798) (6.600,0.1657) (6.900,0.1527) (7.200,0.1407) (7.500,0.1297) (7.800,0.1195) (8.100,0.1101) (8.400,0.1015)};
\foreach \x in {-8,-4,0,4,8} {\node[below] at (\x,0) {$\x$};}
\node[cBar,anchor=west] at (1.6,1.0) {$G(r)=t^{|r|}$，$t=\tanh1=0.7616$};
\node[cExact,anchor=west] at (2.6,0.78) {虛線 $e^{-|r|/\xi}$，$\xi=3.6719$};
\end{tikzpicture}}

\newcommand{\NineFigWell}{%
\begin{tikzpicture}[x=1.35cm,y=1.4cm,>=Stealth,font=\scriptsize,baseline=(current bounding box.center)]
\draw[->] (-2.15,0) -- (2.35,0) node[right] {$m$};
\draw[->] (0,-0.6) -- (0,1.68) node[above] {$F(m)$};
\foreach \x in {-2,-1,1,2} {\draw (\x,0) -- (\x,-0.04);}
\node[below] at (-2,-0.02) {$-2$}; \node[below] at (2,-0.02) {$2$};
\foreach \y in {1} {\draw (0,\y) -- (-0.05,\y) node[left] {$\y$};}
\begin{scope}
\clip (-2.1,-0.6) rectangle (2.1,1.5);
\draw[cBar!35,thick] plot coordinates {(-1.450,2.1564) (-1.400,1.9404) (-1.350,1.7416) (-1.300,1.5590) (-1.250,1.3916) (-1.200,1.2384) (-1.150,1.0985) (-1.100,0.9710) (-1.050,0.8551) (-1.000,0.7500) (-0.950,0.6549) (-0.900,0.5690) (-0.850,0.4918) (-0.800,0.4224) (-0.750,0.3604) (-0.700,0.3050) (-0.650,0.2559) (-0.600,0.2124) (-0.550,0.1741) (-0.500,0.1406) (-0.450,0.1115) (-0.400,0.0864) (-0.350,0.0650) (-0.300,0.0470) (-0.250,0.0322) (-0.200,0.0204) (-0.150,0.0114) (-0.100,0.0050) (-0.050,0.0013) (0.000,0.0000) (0.050,0.0013) (0.100,0.0050) (0.150,0.0114) (0.200,0.0204) (0.250,0.0322) (0.300,0.0470) (0.350,0.0650) (0.400,0.0864) (0.450,0.1115) (0.500,0.1406) (0.550,0.1741) (0.600,0.2124) (0.650,0.2559) (0.700,0.3050) (0.750,0.3604) (0.800,0.4224) (0.850,0.4918) (0.900,0.5690) (0.950,0.6549) (1.000,0.7500) (1.050,0.8551) (1.100,0.9710) (1.150,1.0985) (1.200,1.2384) (1.250,1.3916) (1.300,1.5590) (1.350,1.7416) (1.400,1.9404) (1.450,2.1564)};
\draw[cBar!65,thick] plot coordinates {(-1.700,2.0880) (-1.650,1.8530) (-1.600,1.6384) (-1.550,1.4430) (-1.500,1.2656) (-1.450,1.1051) (-1.400,0.9604) (-1.350,0.8304) (-1.300,0.7140) (-1.250,0.6104) (-1.200,0.5184) (-1.150,0.4373) (-1.100,0.3660) (-1.050,0.3039) (-1.000,0.2500) (-0.950,0.2036) (-0.900,0.1640) (-0.850,0.1305) (-0.800,0.1024) (-0.750,0.0791) (-0.700,0.0600) (-0.650,0.0446) (-0.600,0.0324) (-0.550,0.0229) (-0.500,0.0156) (-0.450,0.0103) (-0.400,0.0064) (-0.350,0.0038) (-0.300,0.0020) (-0.250,0.0010) (-0.200,0.0004) (-0.150,0.0001) (-0.100,0.0000) (-0.050,0.0000) (0.000,0.0000) (0.050,0.0000) (0.100,0.0000) (0.150,0.0001) (0.200,0.0004) (0.250,0.0010) (0.300,0.0020) (0.350,0.0038) (0.400,0.0064) (0.450,0.0103) (0.500,0.0156) (0.550,0.0229) (0.600,0.0324) (0.650,0.0446) (0.700,0.0600) (0.750,0.0791) (0.800,0.1024) (0.850,0.1305) (0.900,0.1640) (0.950,0.2036) (1.000,0.2500) (1.050,0.3039) (1.100,0.3660) (1.150,0.4373) (1.200,0.5184) (1.250,0.6104) (1.300,0.7140) (1.350,0.8304) (1.400,0.9604) (1.450,1.1051) (1.500,1.2656) (1.550,1.4430) (1.600,1.6384) (1.650,1.8530) (1.700,2.0880)};
\draw[cBar,thick] plot coordinates {(-2.000,2.0000) (-1.950,1.7135) (-1.900,1.4530) (-1.850,1.2171) (-1.800,1.0044) (-1.750,0.8135) (-1.700,0.6430) (-1.650,0.4918) (-1.600,0.3584) (-1.550,0.2418) (-1.500,0.1406) (-1.450,0.0539) (-1.400,-0.0196) (-1.350,-0.0809) (-1.300,-0.1310) (-1.250,-0.1709) (-1.200,-0.2016) (-1.150,-0.2240) (-1.100,-0.2390) (-1.050,-0.2474) (-1.000,-0.2500) (-0.950,-0.2476) (-0.900,-0.2410) (-0.850,-0.2307) (-0.800,-0.2176) (-0.750,-0.2021) (-0.700,-0.1850) (-0.650,-0.1666) (-0.600,-0.1476) (-0.550,-0.1284) (-0.500,-0.1094) (-0.450,-0.0910) (-0.400,-0.0736) (-0.350,-0.0575) (-0.300,-0.0430) (-0.250,-0.0303) (-0.200,-0.0196) (-0.150,-0.0111) (-0.100,-0.0050) (-0.050,-0.0012) (0.000,0.0000) (0.050,-0.0012) (0.100,-0.0050) (0.150,-0.0111) (0.200,-0.0196) (0.250,-0.0303) (0.300,-0.0430) (0.350,-0.0575) (0.400,-0.0736) (0.450,-0.0910) (0.500,-0.1094) (0.550,-0.1284) (0.600,-0.1476) (0.650,-0.1666) (0.700,-0.1850) (0.750,-0.2021) (0.800,-0.2176) (0.850,-0.2307) (0.900,-0.2410) (0.950,-0.2476) (1.000,-0.2500) (1.050,-0.2474) (1.100,-0.2390) (1.150,-0.2240) (1.200,-0.2016) (1.250,-0.1709) (1.300,-0.1310) (1.350,-0.0809) (1.400,-0.0196) (1.450,0.0539) (1.500,0.1406) (1.550,0.2418) (1.600,0.3584) (1.650,0.4918) (1.700,0.6430) (1.750,0.8135) (1.800,1.0044) (1.850,1.2171) (1.900,1.4530) (1.950,1.7135) (2.000,2.0000)};
\draw[cWalk,densely dashed,thick] plot coordinates {(-1.950,1.9085) (-1.900,1.6430) (-1.850,1.4021) (-1.800,1.1844) (-1.750,0.9885) (-1.700,0.8130) (-1.650,0.6568) (-1.600,0.5184) (-1.550,0.3968) (-1.500,0.2906) (-1.450,0.1989) (-1.400,0.1204) (-1.350,0.0541) (-1.300,-0.0010) (-1.250,-0.0459) (-1.200,-0.0816) (-1.150,-0.1090) (-1.100,-0.1290) (-1.050,-0.1424) (-1.000,-0.1500) (-0.950,-0.1526) (-0.900,-0.1510) (-0.850,-0.1457) (-0.800,-0.1376) (-0.750,-0.1271) (-0.700,-0.1150) (-0.650,-0.1016) (-0.600,-0.0876) (-0.550,-0.0734) (-0.500,-0.0594) (-0.450,-0.0460) (-0.400,-0.0336) (-0.350,-0.0225) (-0.300,-0.0130) (-0.250,-0.0053) (-0.200,0.0004) (-0.150,0.0039) (-0.100,0.0050) (-0.050,0.0038) (0.000,0.0000) (0.050,-0.0062) (0.100,-0.0150) (0.150,-0.0261) (0.200,-0.0396) (0.250,-0.0553) (0.300,-0.0730) (0.350,-0.0925) (0.400,-0.1136) (0.450,-0.1360) (0.500,-0.1594) (0.550,-0.1834) (0.600,-0.2076) (0.650,-0.2316) (0.700,-0.2550) (0.750,-0.2771) (0.800,-0.2976) (0.850,-0.3157) (0.900,-0.3310) (0.950,-0.3426) (1.000,-0.3500) (1.050,-0.3524) (1.100,-0.3490) (1.150,-0.3390) (1.200,-0.3216) (1.250,-0.2959) (1.300,-0.2610) (1.350,-0.2159) (1.400,-0.1596) (1.450,-0.0911) (1.500,-0.0094) (1.550,0.0868) (1.600,0.1984) (1.650,0.3268) (1.700,0.4730) (1.750,0.6385) (1.800,0.8244) (1.850,1.0321) (1.900,1.2630) (1.950,1.5185) (2.000,1.8000)};
\end{scope}
\fill[cExact] (-1.000,-0.2500) circle (1.6pt); \fill[cExact] (1.000,-0.2500) circle (1.6pt);
\node[cExact,below] at (-1.000,-0.2800) {$m^*=\pm1$，$F_{\min}=-0.25$};
\shade[ball color=cWalk] (1.0467,-0.2824) circle (2.6pt);
\node[cWalk,anchor=west] at (1.147,-0.3724) {$h=0.1$：$m^*=1.0467$};
\node[cBar!50,anchor=south] at (-1.28,1.5) {$r=1$};
\node[cBar!75,anchor=south east] at (1.57,1.5) {$r=0$};
\node[cBar,anchor=south] at (1.91,1.5) {$r=-1$};
\end{tikzpicture}}

\newcommand{\NineTabZmf}{%
\begin{tabular}{@{}lccc@{}}
\toprule
$V$（$r=-1$，$u=1$） & 10 & 100 & 1000\\
\midrule
$\langle m\rangle$（$h=0.01$） & 0.0869 & 0.7576 & 1.0042\\
$\langle|m|\rangle$（$h=0$） & 0.8888 & 0.9922 & 0.9992\\
$-\ln Z/V$（$h=0$） & $-0.2719$ & $-0.2397$ & $-0.2478$\\
\bottomrule
\end{tabular}}

\newcommand{\NineTabExNineFour}{%
\begin{tabular}{@{}l*{5}{c}@{}}
\toprule
練習 9.4 & $\beta$ & $\Delta$ & $\alpha$ & $1/\delta$ & $\gamma$\\
\midrule
數值斜率 & 0.5000 & 1.5000 & 0.0000 & 0.3333 & 1.0000\\
理論 & $\frac{1}{2}$ & $\frac{3}{2}$ & $0$ & $\frac{1}{3}$ & $1$\\
\bottomrule
\end{tabular}}

\newcommand{\NineTabSaw}{%
\begin{tabular}{@{}l*{6}{r}@{}}
\toprule
$N$ & 1 & 2 & 3 & 4 & 5 & 6\\
\midrule
$c_N$ & 4 & 12 & 36 & 100 & 284 & 780\\
$c_N\langle R^2\rangle_N$ & 4 & 32 & 164 & 704 & 2716 & 9808\\
\bottomrule
\end{tabular}}

\newcommand{\NineTabCoef}{%
\begin{tabular}{@{}lccc@{}}
\toprule
精確 $\div$ 漸近，$N=$ & 10 & 100 & 1000\\
\midrule
$s=\gamma=\frac{43}{32}$ & 1.0227 & 1.0023 & 1.0002\\
$s=\gamma+2\nu=\frac{91}{32}$ & 1.2760 & 1.0264 & 1.0026\\
係數比 $\div$ (9.41) & 1.2476 & 1.0240 & 1.0024\\
\bottomrule
\end{tabular}}

\newcommand{\NineFigLadder}{%
\begin{tikzpicture}[font=\scriptsize,baseline=(current bounding box.center)]
\node[anchor=west] at (0,0.21) {$G_{03}$ 的分子};
\draw[gray!35] (1.75,0.00) -- (2.17,0.00);
\draw[gray!35] (2.17,0.00) -- (2.59,0.00);
\draw[gray!35] (1.75,0.42) -- (2.17,0.42);
\draw[gray!35] (2.17,0.42) -- (2.59,0.42);
\draw[gray!35] (1.75,0.00) -- (1.75,0.42);
\draw[gray!35] (2.17,0.00) -- (2.17,0.42);
\draw[gray!35] (2.59,0.00) -- (2.59,0.42);
\draw[cWalk,very thick] (1.75,0.00) -- (1.75,0.42);
\fill[gray!60] (1.75,0.00) circle (0.8pt);
\fill[gray!60] (2.17,0.00) circle (0.8pt);
\fill[gray!60] (2.59,0.00) circle (0.8pt);
\fill[gray!60] (1.75,0.42) circle (0.8pt);
\fill[gray!60] (2.17,0.42) circle (0.8pt);
\fill[gray!60] (2.59,0.42) circle (0.8pt);
\fill[cBar] (1.75,0.00) circle (1.6pt);
\fill[cBar] (1.75,0.42) circle (1.6pt);
\node[below,inner sep=1.5pt] at (2.17,-0.00) {$K$};
\draw[gray!35] (3.05,0.00) -- (3.47,0.00);
\draw[gray!35] (3.47,0.00) -- (3.89,0.00);
\draw[gray!35] (3.05,0.42) -- (3.47,0.42);
\draw[gray!35] (3.47,0.42) -- (3.89,0.42);
\draw[gray!35] (3.05,0.00) -- (3.05,0.42);
\draw[gray!35] (3.47,0.00) -- (3.47,0.42);
\draw[gray!35] (3.89,0.00) -- (3.89,0.42);
\draw[cWalk,very thick] (3.05,0.00) -- (3.47,0.00);
\draw[cWalk,very thick] (3.47,0.00) -- (3.47,0.42);
\draw[cWalk,very thick] (3.05,0.42) -- (3.47,0.42);
\fill[gray!60] (3.05,0.00) circle (0.8pt);
\fill[gray!60] (3.47,0.00) circle (0.8pt);
\fill[gray!60] (3.89,0.00) circle (0.8pt);
\fill[gray!60] (3.05,0.42) circle (0.8pt);
\fill[gray!60] (3.47,0.42) circle (0.8pt);
\fill[gray!60] (3.89,0.42) circle (0.8pt);
\fill[cBar] (3.05,0.00) circle (1.6pt);
\fill[cBar] (3.05,0.42) circle (1.6pt);
\node[below,inner sep=1.5pt] at (3.47,-0.00) {$K^{3}$};
\draw[gray!35] (4.35,0.00) -- (4.77,0.00);
\draw[gray!35] (4.77,0.00) -- (5.19,0.00);
\draw[gray!35] (4.35,0.42) -- (4.77,0.42);
\draw[gray!35] (4.77,0.42) -- (5.19,0.42);
\draw[gray!35] (4.35,0.00) -- (4.35,0.42);
\draw[gray!35] (4.77,0.00) -- (4.77,0.42);
\draw[gray!35] (5.19,0.00) -- (5.19,0.42);
\draw[cWalk,very thick] (4.35,0.00) -- (4.77,0.00);
\draw[cWalk,very thick] (4.77,0.00) -- (5.19,0.00);
\draw[cWalk,very thick] (5.19,0.00) -- (5.19,0.42);
\draw[cWalk,very thick] (4.77,0.42) -- (5.19,0.42);
\draw[cWalk,very thick] (4.35,0.42) -- (4.77,0.42);
\fill[gray!60] (4.35,0.00) circle (0.8pt);
\fill[gray!60] (4.77,0.00) circle (0.8pt);
\fill[gray!60] (5.19,0.00) circle (0.8pt);
\fill[gray!60] (4.35,0.42) circle (0.8pt);
\fill[gray!60] (4.77,0.42) circle (0.8pt);
\fill[gray!60] (5.19,0.42) circle (0.8pt);
\fill[cBar] (4.35,0.00) circle (1.6pt);
\fill[cBar] (4.35,0.42) circle (1.6pt);
\node[below,inner sep=1.5pt] at (4.77,-0.00) {$K^{5}$};
\draw[gray!35] (5.65,0.00) -- (6.07,0.00);
\draw[gray!35] (6.07,0.00) -- (6.49,0.00);
\draw[gray!35] (5.65,0.42) -- (6.07,0.42);
\draw[gray!35] (6.07,0.42) -- (6.49,0.42);
\draw[gray!35] (5.65,0.00) -- (5.65,0.42);
\draw[gray!35] (6.07,0.00) -- (6.07,0.42);
\draw[gray!35] (6.49,0.00) -- (6.49,0.42);
\draw[cWalk,very thick] (5.65,0.00) -- (5.65,0.42);
\draw[cWrong,very thick] (6.07,0.00) -- (6.49,0.00);
\draw[cWrong,very thick] (6.07,0.42) -- (6.49,0.42);
\draw[cWrong,very thick] (6.07,0.00) -- (6.07,0.42);
\draw[cWrong,very thick] (6.49,0.00) -- (6.49,0.42);
\fill[gray!60] (5.65,0.00) circle (0.8pt);
\fill[gray!60] (6.07,0.00) circle (0.8pt);
\fill[gray!60] (6.49,0.00) circle (0.8pt);
\fill[gray!60] (5.65,0.42) circle (0.8pt);
\fill[gray!60] (6.07,0.42) circle (0.8pt);
\fill[gray!60] (6.49,0.42) circle (0.8pt);
\fill[cBar] (5.65,0.00) circle (1.6pt);
\fill[cBar] (5.65,0.42) circle (1.6pt);
\node[below,inner sep=1.5pt] at (6.07,-0.00) {$nK^{5}$};
\node[anchor=west] at (0,-0.87) {$G_{05}$ 的分子};
\draw[gray!35] (1.75,-1.08) -- (2.17,-1.08);
\draw[gray!35] (2.17,-1.08) -- (2.59,-1.08);
\draw[gray!35] (1.75,-0.66) -- (2.17,-0.66);
\draw[gray!35] (2.17,-0.66) -- (2.59,-0.66);
\draw[gray!35] (1.75,-1.08) -- (1.75,-0.66);
\draw[gray!35] (2.17,-1.08) -- (2.17,-0.66);
\draw[gray!35] (2.59,-1.08) -- (2.59,-0.66);
\draw[cWalk,very thick] (1.75,-1.08) -- (2.17,-1.08);
\draw[cWalk,very thick] (2.17,-1.08) -- (2.59,-1.08);
\draw[cWalk,very thick] (2.59,-1.08) -- (2.59,-0.66);
\fill[gray!60] (1.75,-1.08) circle (0.8pt);
\fill[gray!60] (2.17,-1.08) circle (0.8pt);
\fill[gray!60] (2.59,-1.08) circle (0.8pt);
\fill[gray!60] (1.75,-0.66) circle (0.8pt);
\fill[gray!60] (2.17,-0.66) circle (0.8pt);
\fill[gray!60] (2.59,-0.66) circle (0.8pt);
\fill[cBar] (1.75,-1.08) circle (1.6pt);
\fill[cBar] (2.59,-0.66) circle (1.6pt);
\node[below,inner sep=1.5pt] at (2.17,-1.08) {$K^{3}$};
\draw[gray!35] (3.05,-1.08) -- (3.47,-1.08);
\draw[gray!35] (3.47,-1.08) -- (3.89,-1.08);
\draw[gray!35] (3.05,-0.66) -- (3.47,-0.66);
\draw[gray!35] (3.47,-0.66) -- (3.89,-0.66);
\draw[gray!35] (3.05,-1.08) -- (3.05,-0.66);
\draw[gray!35] (3.47,-1.08) -- (3.47,-0.66);
\draw[gray!35] (3.89,-1.08) -- (3.89,-0.66);
\draw[cWalk,very thick] (3.05,-1.08) -- (3.47,-1.08);
\draw[cWalk,very thick] (3.47,-1.08) -- (3.47,-0.66);
\draw[cWalk,very thick] (3.47,-0.66) -- (3.89,-0.66);
\fill[gray!60] (3.05,-1.08) circle (0.8pt);
\fill[gray!60] (3.47,-1.08) circle (0.8pt);
\fill[gray!60] (3.89,-1.08) circle (0.8pt);
\fill[gray!60] (3.05,-0.66) circle (0.8pt);
\fill[gray!60] (3.47,-0.66) circle (0.8pt);
\fill[gray!60] (3.89,-0.66) circle (0.8pt);
\fill[cBar] (3.05,-1.08) circle (1.6pt);
\fill[cBar] (3.89,-0.66) circle (1.6pt);
\node[below,inner sep=1.5pt] at (3.47,-1.08) {$K^{3}$};
\draw[gray!35] (4.35,-1.08) -- (4.77,-1.08);
\draw[gray!35] (4.77,-1.08) -- (5.19,-1.08);
\draw[gray!35] (4.35,-0.66) -- (4.77,-0.66);
\draw[gray!35] (4.77,-0.66) -- (5.19,-0.66);
\draw[gray!35] (4.35,-1.08) -- (4.35,-0.66);
\draw[gray!35] (4.77,-1.08) -- (4.77,-0.66);
\draw[gray!35] (5.19,-1.08) -- (5.19,-0.66);
\draw[cWalk,very thick] (4.35,-1.08) -- (4.35,-0.66);
\draw[cWalk,very thick] (4.35,-0.66) -- (4.77,-0.66);
\draw[cWalk,very thick] (4.77,-0.66) -- (5.19,-0.66);
\fill[gray!60] (4.35,-1.08) circle (0.8pt);
\fill[gray!60] (4.77,-1.08) circle (0.8pt);
\fill[gray!60] (5.19,-1.08) circle (0.8pt);
\fill[gray!60] (4.35,-0.66) circle (0.8pt);
\fill[gray!60] (4.77,-0.66) circle (0.8pt);
\fill[gray!60] (5.19,-0.66) circle (0.8pt);
\fill[cBar] (4.35,-1.08) circle (1.6pt);
\fill[cBar] (5.19,-0.66) circle (1.6pt);
\node[below,inner sep=1.5pt] at (4.77,-1.08) {$K^{3}$};
\draw[gray!35] (5.65,-1.08) -- (6.07,-1.08);
\draw[gray!35] (6.07,-1.08) -- (6.49,-1.08);
\draw[gray!35] (5.65,-0.66) -- (6.07,-0.66);
\draw[gray!35] (6.07,-0.66) -- (6.49,-0.66);
\draw[gray!35] (5.65,-1.08) -- (5.65,-0.66);
\draw[gray!35] (6.07,-1.08) -- (6.07,-0.66);
\draw[gray!35] (6.49,-1.08) -- (6.49,-0.66);
\draw[cWalk,very thick] (5.65,-1.08) -- (5.65,-0.66);
\draw[cWalk,very thick] (5.65,-0.66) -- (6.07,-0.66);
\draw[cWalk,very thick] (6.07,-1.08) -- (6.07,-0.66);
\draw[cWalk,very thick] (6.07,-1.08) -- (6.49,-1.08);
\draw[cWalk,very thick] (6.49,-1.08) -- (6.49,-0.66);
\fill[gray!60] (5.65,-1.08) circle (0.8pt);
\fill[gray!60] (6.07,-1.08) circle (0.8pt);
\fill[gray!60] (6.49,-1.08) circle (0.8pt);
\fill[gray!60] (5.65,-0.66) circle (0.8pt);
\fill[gray!60] (6.07,-0.66) circle (0.8pt);
\fill[gray!60] (6.49,-0.66) circle (0.8pt);
\fill[cBar] (5.65,-1.08) circle (1.6pt);
\fill[cBar] (6.49,-0.66) circle (1.6pt);
\node[below,inner sep=1.5pt] at (6.07,-1.08) {$K^{5}$};
\node[anchor=west] at (0,-1.95) {分母 $Z$};
\draw[gray!35] (1.75,-2.16) -- (2.17,-2.16);
\draw[gray!35] (2.17,-2.16) -- (2.59,-2.16);
\draw[gray!35] (1.75,-1.74) -- (2.17,-1.74);
\draw[gray!35] (2.17,-1.74) -- (2.59,-1.74);
\draw[gray!35] (1.75,-2.16) -- (1.75,-1.74);
\draw[gray!35] (2.17,-2.16) -- (2.17,-1.74);
\draw[gray!35] (2.59,-2.16) -- (2.59,-1.74);
\fill[gray!60] (1.75,-2.16) circle (0.8pt);
\fill[gray!60] (2.17,-2.16) circle (0.8pt);
\fill[gray!60] (2.59,-2.16) circle (0.8pt);
\fill[gray!60] (1.75,-1.74) circle (0.8pt);
\fill[gray!60] (2.17,-1.74) circle (0.8pt);
\fill[gray!60] (2.59,-1.74) circle (0.8pt);
\node[below,inner sep=1.5pt] at (2.17,-2.16) {$1$};
\draw[gray!35] (3.05,-2.16) -- (3.47,-2.16);
\draw[gray!35] (3.47,-2.16) -- (3.89,-2.16);
\draw[gray!35] (3.05,-1.74) -- (3.47,-1.74);
\draw[gray!35] (3.47,-1.74) -- (3.89,-1.74);
\draw[gray!35] (3.05,-2.16) -- (3.05,-1.74);
\draw[gray!35] (3.47,-2.16) -- (3.47,-1.74);
\draw[gray!35] (3.89,-2.16) -- (3.89,-1.74);
\draw[cWrong,very thick] (3.05,-2.16) -- (3.47,-2.16);
\draw[cWrong,very thick] (3.05,-1.74) -- (3.47,-1.74);
\draw[cWrong,very thick] (3.05,-2.16) -- (3.05,-1.74);
\draw[cWrong,very thick] (3.47,-2.16) -- (3.47,-1.74);
\fill[gray!60] (3.05,-2.16) circle (0.8pt);
\fill[gray!60] (3.47,-2.16) circle (0.8pt);
\fill[gray!60] (3.89,-2.16) circle (0.8pt);
\fill[gray!60] (3.05,-1.74) circle (0.8pt);
\fill[gray!60] (3.47,-1.74) circle (0.8pt);
\fill[gray!60] (3.89,-1.74) circle (0.8pt);
\node[below,inner sep=1.5pt] at (3.47,-2.16) {$nK^{4}$};
\draw[gray!35] (4.35,-2.16) -- (4.77,-2.16);
\draw[gray!35] (4.77,-2.16) -- (5.19,-2.16);
\draw[gray!35] (4.35,-1.74) -- (4.77,-1.74);
\draw[gray!35] (4.77,-1.74) -- (5.19,-1.74);
\draw[gray!35] (4.35,-2.16) -- (4.35,-1.74);
\draw[gray!35] (4.77,-2.16) -- (4.77,-1.74);
\draw[gray!35] (5.19,-2.16) -- (5.19,-1.74);
\draw[cWrong,very thick] (4.77,-2.16) -- (5.19,-2.16);
\draw[cWrong,very thick] (4.77,-1.74) -- (5.19,-1.74);
\draw[cWrong,very thick] (4.77,-2.16) -- (4.77,-1.74);
\draw[cWrong,very thick] (5.19,-2.16) -- (5.19,-1.74);
\fill[gray!60] (4.35,-2.16) circle (0.8pt);
\fill[gray!60] (4.77,-2.16) circle (0.8pt);
\fill[gray!60] (5.19,-2.16) circle (0.8pt);
\fill[gray!60] (4.35,-1.74) circle (0.8pt);
\fill[gray!60] (4.77,-1.74) circle (0.8pt);
\fill[gray!60] (5.19,-1.74) circle (0.8pt);
\node[below,inner sep=1.5pt] at (4.77,-2.16) {$nK^{4}$};
\draw[gray!35] (5.65,-2.16) -- (6.07,-2.16);
\draw[gray!35] (6.07,-2.16) -- (6.49,-2.16);
\draw[gray!35] (5.65,-1.74) -- (6.07,-1.74);
\draw[gray!35] (6.07,-1.74) -- (6.49,-1.74);
\draw[gray!35] (5.65,-2.16) -- (5.65,-1.74);
\draw[gray!35] (6.07,-2.16) -- (6.07,-1.74);
\draw[gray!35] (6.49,-2.16) -- (6.49,-1.74);
\draw[cWrong,very thick] (5.65,-2.16) -- (6.07,-2.16);
\draw[cWrong,very thick] (6.07,-2.16) -- (6.49,-2.16);
\draw[cWrong,very thick] (5.65,-1.74) -- (6.07,-1.74);
\draw[cWrong,very thick] (6.07,-1.74) -- (6.49,-1.74);
\draw[cWrong,very thick] (5.65,-2.16) -- (5.65,-1.74);
\draw[cWrong,very thick] (6.49,-2.16) -- (6.49,-1.74);
\fill[gray!60] (5.65,-2.16) circle (0.8pt);
\fill[gray!60] (6.07,-2.16) circle (0.8pt);
\fill[gray!60] (6.49,-2.16) circle (0.8pt);
\fill[gray!60] (5.65,-1.74) circle (0.8pt);
\fill[gray!60] (6.07,-1.74) circle (0.8pt);
\fill[gray!60] (6.49,-1.74) circle (0.8pt);
\node[below,inner sep=1.5pt] at (6.07,-2.16) {$nK^{6}$};
\end{tikzpicture}}

\newcommand{\NineTabLadder}{%
\begin{tabular}{@{}lcccc@{}}
\toprule
$K=0.5$ & $n=1$ & $n=2$ & $n\to0$ & 自迴避路徑\\
\midrule
$G_{03}$ & 0.6027 & 0.5610 & 0.65625 & $K+K^3+K^5$\\
$G_{05}$ & 0.3562 & 0.3171 & 0.40625 & $3K^3+K^5$\\
\bottomrule
\end{tabular}}

\newcommand{\NineTabTrotter}{%
\begin{tabular}{@{}lccc@{}}
\toprule
$\epsilon$ & 0.1 & 0.05 & 0.025\\
\midrule
一階誤差 & $0.005784$ & $0.001502$ & $0.0003827$\\
乘上 $e^{\frac{i\epsilon^2}4[\phi,\nabla^2]}$ & $0.0009465$ & $0.000124$ & $1.589\times10^{-5}$\\
\midrule
$t=1$ 切成 $M$ 段 & $M=10$ & $M=20$ & $M=40$\\
\quad 總誤差 & $0.02366$ & $0.0117$ & $0.005821$\\
\bottomrule
\end{tabular}}

\newcommand{\NineTabWeak}{%
\begin{tabular}{@{}l*{4}{r}r@{}}
\toprule
$u$ & 0 & 1 & 4 & $\infty$ & 書上 $u=1$\\
\midrule
$\sum_we^{-(u/2)n_w}$ & 65536 & 23695.11 & 7988.15 & 5916 & {\color{cWrong}$6.2\times10^{5}$}\\
\bottomrule
\end{tabular}}

\newcommand{\NineTabErrata}{%
\begin{tabular}{@{}lll@{}}
\toprule
位置 & 書上 & 應為\\
\midrule
p.~210，Table 9.1 三維欄 & {\color{cWrong}$\alpha=\alpha'=0.12$} & {\color{cExact}$\approx0.19$（(9.16)(9.18) 都得 0.194）}\\
p.~210，Table 9.1 第二個關聯長度 & {\color{cWrong}$\xi=\int C\rho^2d^d\rho/\!\int C\,d^d\rho$} & {\color{cExact}$\xi^2=\int C\rho^2d^d\rho/\!\int C\,d^d\rho$}\\
p.~218，(9.37) 之後 & {\color{cWrong}關聯長度 $\propto K_1/(z_c-z)^\nu$} & {\color{cExact}$\xi^2\approx K_1/(z_c-z)^{2\nu}$}\\
p.~218，(9.38)$\to$(9.39) & {\color{cWrong}$z_c^{-(N+s)}N^{\delta-1}/\Gamma(s)$} & {\color{cExact}$z_c^{-(N+s)}N^{s-1}/\Gamma(s)$}\\
p.~219，(9.40)$\to$(9.41) & {\color{cWrong}must be equal to $V\Gamma_N/z_0$} & {\color{cExact}$V\Gamma_N/z_0^N$}\\
p.~219，(9.41) & {\color{cWrong}$K_1z_c^{-\gamma}\,\Gamma(\gamma)N^{2\nu}/\Gamma(\gamma+2\nu)$} & {\color{cExact}$K_1z_c^{-2\nu}\,\Gamma(\gamma)N^{2\nu}/\Gamma(\gamma+2\nu)$}\\
p.~225，(9.70) & {\color{cWrong}$e^{\nabla^2\epsilon}$ 的核 $(2\pi)^{-d/2}e^{-\frac d2|\Delta\vec R|^2/\epsilon}$} & {\color{cExact}$e^{\frac12\nabla^2\epsilon}$ 的核 $(2\pi\epsilon)^{-d/2}e^{-|\Delta\vec R|^2/2\epsilon}$}\\
p.~226，(9.72) & {\color{cWrong}$\int_0^\infty$} & {\color{cExact}$\int_{-\infty}^{\infty}$}\\
p.~226，(9.75) & {\color{cWrong}$-\int\frac d2\big|\frac{d\vec R}{ds}\big|^2ds\ \frac u4\iint\delta$} & {\color{cExact}$-\int\frac d2\big|\frac{d\vec R}{ds}\big|^2ds-\frac u4\iint\delta$}\\
p.~230，(9:S-19) & {\color{cWrong}$\partial^3/\partial x_r\partial x_s\partial x_t$；括號裡的 $\langle x_rx_s\rangle$} & {\color{cExact}$\partial^3/\partial h_r\partial h_s\partial h_t$；$\langle x_tx_u\rangle$}\\
\bottomrule
\end{tabular}}


\section{漫步的標度與臨界現象（第 9 章）}

\begin{keybox}[先看這裡]
第 7、8 章得到自迴避漫步 $\langle R^2\rangle\propto N^{2\nu}$。這一章說明 de Gennes 的發現：\textbf{$\nu$ 是一個磁鐵的臨界指數}——「每個自旋有 $n$ 個分量」的磁鐵
（O($n$) 模型）在 $n\to0$ 時，關聯函數就是自迴避漫步的生成函數。路線：統計力學的詞彙（\textbf{配分函數、磁化率、關聯長度、臨界指數}）$\to$ 六個指數之間的\textbf{標度律}
$\to$ 讓自旋可以伸縮的 Ginzburg--Landau--Wilson 模型 $\to$ 取 $z^N$ 的係數得 $\langle R^2\rangle\propto N^{2\nu}$ $\to$ 用高斯積分說明 $n\to0$ 時只剩不碰自己的路徑。
不熟統計力學也能讀：每個概念都先用 2--16 個自旋的小例子算到底。

\textbf{記號}：$J=k_B=1$；Ising 自旋 $s=\pm1$、$K=J/T$、無因次外場 $H=h/T$、每個自旋的磁化 $m$；O($n$) 的分量數 $n$；臨界指數 $\alpha,\beta,\gamma,\delta,\nu,\eta$；
GLW 的約化溫度 $r$（$\propto T-T_c$）與四次項係數 $u$；自迴避條數 $c_N$（書上 9.4 節叫 $\Gamma_N$）、每步的逸度 $z$、$z_c=1/\mu$（二維 0.3791）。

\textbf{全章的主例}：4 個自旋 $Z=2e^{4K}+12+2e^{-4K}$；一維 Ising $\chi=e^{2K}$；雙井 $m^*=\sqrt{-r/u}$；梯子在 $n\to0$ 時 $G_{03}=K+K^3+K^5$（三條自迴避路徑）。
\begin{center}
\small
\begin{tabular}{@{}lll@{\qquad}lll@{}}
\toprule
節 & 內容 & 動畫 & 節 & 內容 & 動畫\\
\midrule
9.1--9.2 & 漫步的指數、臨界點 & W901 & 9.3 & Ginzburg--Landau--Wilson、雙井 & W904\\
9.2.1--9.2.2 & O($n$) 模型、對稱破缺 & W902 & 9.4 & 從磁鐵讀出 $\langle R^2\rangle$ & W905\\
9.2.2--9.2.3 & 臨界指數、標度律 & W903 & 9.5--9.6 & $n\to0$、路徑積分、高斯積分 & W906\\
\bottomrule
\end{tabular}
\end{center}
（上表是書上的節號，本講義 9.1--9.6 依序對應，書上的高斯積分補充在 9.6 節末。）
數字一律由 \texttt{codes/utils/walk\_crit.py} 重算，逐式校對見 \texttt{ref/md/ch09\_formulas.md}，圖表由 \texttt{ch9/make\_figs.py} 產生並先做數字檢查。連續播放：\texttt{Ch9Demo}（W901--W906）。
\end{keybox}

% =====================================================================
\subsection{漫步的指數就是臨界指數}
\leavevmode\bk{pp.~203--205}\ \anim{W901}\par\nopagebreak

\textbf{問}：自迴避漫步越走越大、磁鐵加熱就失去磁性——為什麼說這是同一個數學問題？

\textbf{漫步的指數}（(9.1)(9.2)）：$\langle R^2\rangle\propto N^p$，Flory（第 7、8 章兩次推出）給 $p=2\nu=\frac6{d+2}$。$d=1,2,3,4$：$p=2,\frac32,\frac65,1$；已知的 $2\nu$ 是 $2,\frac32,1.176,1$（圖 \ref{nine-fig-ms} 右上）：一、二、四維完全正確，三維 0.600 對 0.588。
〔注〕書上說三維「差一個百分點左右」；實際 $0.600/0.588-1=0.0204$，差 2\%。二維精確列舉（第 7 章）的局部指數 $2\nu_{\rm eff}=1.4442,\ 1.4505,\ 1.4556$（$N=10,12,14$），慢慢爬向 $\frac32$。

\textbf{磁鐵的三分鐘入門}：鐵塊裡每個原子像一根小磁針（\textbf{自旋}）。相鄰自旋同向時能量較低，所以低溫時它們排成一致的方向，沒有外加磁場也有\textbf{自發磁化} $M_s$；
溫度把自旋打亂，高過\textbf{臨界溫度} $T_c$（居禮溫度）就 $M_s=0$。重點是 $T_c$ 附近 $M_s$ \emph{怎麼}消失：
\[M_s\propto(T_c-T)^\beta,\tag{9.3}\]
$\beta$ 叫\textbf{臨界指數}。它和漫步的 $\frac12$ 一樣是\textbf{普適}的：鐵、鎳、二元合金……細節不同，只要空間維度與「自旋有幾個分量」相同，$\beta$ 就相同。

\textbf{兩個算到底的例子}（圖 \ref{nine-fig-ms} 左）。(1) \textbf{平均場}：把一個自旋的鄰居都換成平均值 $m$，它在這個「平均的場」裡的平均是 $\tanh(m/T)$（$T$ 以 $T_c$ 為單位），自洽要求
$m=\tanh(m/T)$。$T_c=1$ 附近用 $\tanh x\approx x-\frac{x^3}3$：除以 $m$ 得 $1=\frac1T-\frac{m^2}{3T^3}$，$m^2=3T^2(1-T)$，$m\approx\sqrt{3(1-T)}$，所以 $\beta=\frac12$
（$T=0.99$：近似式 0.1732、精確解 0.1725）。(2) \textbf{二維 Ising}（Onsager 的精確解）：$M_s=\big(1-\sinh^{-4}(2/T)\big)^{1/8}$，$T_c=2/\ln(1+\sqrt2)=2.2692$（$\sinh(2/T_c)=1$），$\beta=\frac18$。
$T/T_c=0.5,\ 0.9,\ 0.99$：平均場 0.9575、0.5254、0.1725；Onsager 0.9980、0.8959、0.6859——$\beta$ 越小，在 $T_c$ 前掉得越陡。$\ln M_s$ 對 $\ln(1-T/T_c)$ 的斜率量到 0.5000、0.1250。
〔注〕書上說 $\beta\approx\frac13$：那是三維（三維 Ising 約 0.326、三維 $n=0$ 約 0.302，只引用）。

\textbf{de Gennes 的字典}：$p=2\nu$ 的 $\nu$，是一個假想磁鐵（$n=0$ 個分量的 O($n$) 模型）的\textbf{關聯長度指數}。對照（9.5、9.6 節證明）：
步數 $N\leftrightarrow1/(z_c-z)$（離臨界點多遠）；$\langle R^2\rangle\propto N^{2\nu}\leftrightarrow\xi^2\propto(z_c-z)^{-2\nu}$；條數 $c_N\propto\mu^NN^{\gamma-1}\leftrightarrow\chi\propto(z_c-z)^{-\gamma}$；不碰自己 $\leftrightarrow n=0$。

\begin{figure}[htbp]
\centering
\begin{minipage}[c]{0.52\linewidth}\centering\NineFigMs\end{minipage}\hfill
\begin{minipage}[c]{0.46\linewidth}\centering\footnotesize\setlength{\tabcolsep}{3.5pt}\NineTabFlory\\[8pt]\NineTabMs\end{minipage}
\caption{漫步的指數與自發磁化。左：$M_s$ 對 $T/T_c$，平均場（橙）$\propto(1-T/T_c)^{1/2}$、二維 Ising（藍）$\propto(1-T/T_c)^{1/8}$；圓點是右下表的三個溫度。
右上：Flory 的 $p$ 與已知的 $2\nu$（三維是數值，只引用）。右下：$M_s$ 與 log--log 斜率}\label{nine-fig-ms}
\end{figure}

\textbf{揭曉}：漫步的 $\nu$ 和磁鐵的 $\beta$ 都是臨界指數；本章要說明為什麼「$n=0$ 的磁鐵」會數自迴避漫步。

% =====================================================================
\subsection{\texorpdfstring{O($n$)}{O(n)} 模型與對稱破缺}
\leavevmode\bk{pp.~205--211}\ \anim{W902}\par\nopagebreak

\textbf{問}：能量在「所有自旋一起轉」之下不變，所以平均磁化一定是 0——那永久磁鐵是怎麼來的？

\textbf{O($n$) 模型}（(9.4)(9.5)）：每個格點一個 $n$ 分量的向量 $\vec S_i$；$n=1$ 是 Ising（只能上下）、$n=2$ 是 xy（在平面上轉；超導、超流）、$n=3$ 是 Heisenberg（空間任意方向），$n=0$ 是自迴避漫步（9.6 節）：
\[H_{\rm int}=-J\sum_{\langle i,j\rangle}\vec S_i\cdot\vec S_j,\qquad\vec M=\sum_i\vec S_i .\tag{9.4, 9.5}\]
只用到內積，所以\textbf{全部一起轉}時 $H$ 不變。例：8 個平面自旋（種子 902）全部轉 $120^\circ$，能量 $-6.3851$、$|\vec M|=6.7943$ 都不變，方向從 $-13.1^\circ$ 變成 $106.9^\circ$。
各個方向一樣好，所以 $\langle\vec M\rangle=0$。

\textbf{統計力學的四個詞}（(9.6)--(9.13)）：加上外場，$H=-J\sum\vec S_i\cdot\vec S_j-\vec h\cdot\sum_i\vec S_i$。溫度 $T$ 時，每個組態出現的機率正比於\textbf{波茲曼權重} $e^{-H/k_BT}$：
\[Z=\sum_{\text{組態}}e^{-H/k_BT},\qquad F=-k_BT\ln Z,\qquad\langle\vec M\rangle=-\frac{\partial F}{\partial\vec h},\qquad\chi_T=\frac{\partial\langle\vec M\rangle}{\partial\vec h}\propto|T_c-T|^{-\gamma}.\tag{9.8--9.13}\]
\textbf{配分函數} $Z$ 是權重的總和（機率的分母），\textbf{自由能} $F$ 是它的對數。$Z$ 就像第 1 章的生成函數：對 $h$ 微分一次得\textbf{平均磁化}，再微分一次得\textbf{磁化率} $\chi$
——磁化對外場有多敏感，它也等於磁化的漲落 $(\langle M^2\rangle-\langle M\rangle^2)/k_BT$。

\textbf{算到底：4 個自旋圍一圈}（$2\times2$ 方塊、4 條鍵、16 個組態）：全同向 2 個（$E=-4$，權重 $e^{4K}$）；翻一個 8 個（$E=0$）；兩上兩下相鄰 4 個（$E=0$）、交錯 2 個（$E=+4$，權重 $e^{-4K}$）。所以
$Z=2e^{4K}+12+2e^{-4K}$，$P(M=\pm4)=e^{4K}/Z$、$P(\pm2)=4/Z$、$P(0)=(4+2e^{-4K})/Z$。$K=1$（$T=1$）：$Z=121.2329$，$P(M=\pm4,\pm2,0)=0.4504,\ 0.0330,\ 0.0333$；$K=0.2$（$T=5$）：$Z=17.3497$，$P=0.1283,\ 0.2306,\ 0.2823$（圖 \ref{nine-fig-break} 左上）。
兩種溫度都左右對稱，$\langle M\rangle=0$；可是 $\langle|m|\rangle=0.9337$ 對 0.4871：低溫時磁化「不是正就是負」，只是正負各半。（2 個自旋、$K=1$：$Z=2(e+e^{-1})=6.1723$、$\langle|m|\rangle=0.8808$。）

\textbf{要等多久才翻面？}用 Metropolis 模擬：每次試翻一個自旋，能量降低就翻，升高 $\Delta E$ 就以機率 $e^{-\Delta E/T}$ 翻；這樣產生的組態，機率正比於 $e^{-E/T}$。
二維 $L\times L$ 週期格子，$T=2.0<T_c$（$M_s=0.9113$），從全部朝上開始，4000 sweep $\times$ 50 個複本（種子 902），$m$ 在 $>0.5$ 與 $<-0.5$ 之間換邊算一次翻轉：
平均等待 133、608、2151 sweep（$L=4,6,8$）——$L$ 每加 2 約乘 4（圖 \ref{nine-fig-break} 右）。取樣器的核對：$L=4$ 的 $\langle|m|\rangle=0.9182$（精確列舉 0.9189）、每個自旋的能量 $-1.7528$（$-1.7554$）。
〔注〕書上說等待時間是「天文數字」，而且系統越大越久。無限大的系統永遠等不到：它選定一邊就不回頭——這就是\textbf{對稱破缺}。

\textbf{外場與相圖}（Fig.~9.2）：加一個很小的外場（$K=1$）：$H=0.1,\ 0.01$ 時 4 個自旋的 $\langle m\rangle=0.3506,\ 0.0367$，$4\times4$ 週期格子 0.9210、0.1584——格子越大，同一個小外場越接近無限大格子的 $M_s=0.9993$。
正確的順序是\textbf{先}讓系統變成無限大、\textbf{再}讓 $h\to0$，得 $M\to\pm M_s$。所以 $T<T_c$、$h=0$ 是一條\textbf{一級相變線}（$h$ 由負變正，$M$ 從 $-M_s$ 跳到 $+M_s$），終點 $(T_c,0)$ 是\textbf{臨界點}。

\textbf{關聯函數與關聯長度}（(9.14)(9.15)）：相距 $\rho$ 的兩個自旋同向的程度
\[C(\vec\rho)=\langle\vec S(\vec R)\cdot\vec S(\vec R+\vec\rho)\rangle\propto\frac{f(|\vec\rho|/\xi)}{|\vec\rho|^{d-2+\eta}},\qquad\xi\propto|T-T_c|^{-\nu}.\tag{9.14, 9.15}\]
$\xi$ 是「同向小島」的大小；靠近 $T_c$ 時各種大小的小島都有，$\xi\to\infty$。

\begin{figure}[htbp]
\centering
\begin{minipage}[c]{0.44\linewidth}\centering\NineFigPM\\[8pt]\footnotesize\setlength{\tabcolsep}{3pt}\NineTabField\end{minipage}\hfill
\begin{minipage}[c]{0.54\linewidth}\centering\NineFigTrace\\[6pt]\footnotesize\setlength{\tabcolsep}{4pt}\NineTabFlip\end{minipage}
\caption{對稱破缺。左上：4 個自旋一圈的 $P(M)$，低溫（藍）集中在 $\pm4$、高溫（橙）集中在 0，都左右對稱。左下：小外場下的 $\langle m\rangle$。
右：Metropolis 的 $m$（每條 200 點），低溫時格子越大越難翻面；下表是 50 個複本的平均翻轉等待}\label{nine-fig-break}
\end{figure}

\textbf{揭曉}：對稱說平均是 0；可是低溫時翻一次要等很久——格子越大越久，無限大就永遠不翻：這就是對稱破缺。

% =====================================================================
\subsection{臨界指數與標度律}
\leavevmode\bk{pp.~209--215}\ \anim{W903}\par\nopagebreak

\textbf{問}：一個磁鐵有六個臨界指數 $\alpha,\beta,\gamma,\delta,\nu,\eta$——它們彼此獨立嗎？

\textbf{六個指數}（Table 9.1，圖 \ref{nine-fig-exp} 左上）：比熱 $\propto|T-T_c|^{-\alpha}$、自發磁化 $\propto(T_c-T)^\beta$、磁化率 $\propto|T-T_c|^{-\gamma}$、$T=T_c$ 時 $M\propto|h|^{1/\delta}$、
關聯長度 $\xi\propto|T-T_c|^{-\nu}$、$T=T_c$ 時 $C(\rho)\propto\rho^{-(d-2+\eta)}$。
\begin{errbox}[p.~210，Table 9.1]
(1) 三維欄的 $\alpha=\alpha'=0.12$ 應約為 0.19：同一欄的 $\beta,\gamma$ 代 (9.16) 得 $2-2\beta-\gamma=0.194$，$\nu$ 代 (9.18) 得 $2-3\nu=0.194$（照書上 $\alpha+2\beta+\gamma=1.926\ne2$；同欄其他值彼此自洽，見圖 \ref{nine-fig-exp} 右）。(2) 第二個「關聯長度」列把定義寫成 $\xi=\int C\rho^2d^d\rho/\!\int C\,d^d\rho$，應為 $\xi^2$（第一列寫對了）：比值是長度的\emph{平方}，見下面一維的例子。
\end{errbox}
〔注〕表中的「$\xi^2\propto|T-T_c|^{-\nu}$」要讀成 $\xi\propto|T-T_c|^{-\nu}$。三維欄是較舊的估計；目前的文獻值（只引用）$\nu\approx0.5876$、$\gamma\approx1.157$、$\eta\approx0.031$、$\alpha\approx0.237$、$\beta\approx0.302$。

\textbf{算到底：一維 Ising 鏈}（$K=1$）：每條鍵的權重寫成 $e^{Kss'}=\cosh K\,(1+t\,ss')$，$t=\tanh K=0.7616$。乘開之後對自旋求和，凡是有自旋只出現一次的項都是 0，
只剩「把 0 和 $r$ 連起來」的那一串鍵（9.6 節的梯子用的是同一招）：
\[G(r)=\langle s_0s_r\rangle=t^{|r|}=1,\ 0.7616,\ 0.5800,\ 0.4417,\ 0.3364,\ \dots,\qquad\xi=-\frac1{\ln t}=3.6719\]
（$G=e^{-|r|/\xi}$，圖 \ref{nine-fig-exp} 左下）。磁化率是關聯的總和（(9.24)--(9.26)：$k_BT\chi_T=\sum_jC(\vec R_i-\vec R_j)$），這裡 $\chi=\sum_rG(r)=1+\frac{2t}{1-t}=\frac{1+t}{1-t}=e^{2K}=7.3891$；
一根一根加，$|r|\le5,10,20$：5.7520、6.9696、7.3615；16 個自旋的環逐一列舉得 7.2021（有限大小）。
Table 9.1 的比值 $\sum r^2G/\sum G=2t/(1-t)^2$ 除以 $\xi^2$ 得 1.9877、1.9998、2.0000（$K=1,2,3$），是常數；除以 $\xi$ 則一直變大——所以它是 $\xi^2$。

\textbf{Fisher 定律的由來}（(9.26)(9.27)）：把 (9.14) 代進 $\chi\propto\int d^dR\,C$，令 $\vec R=\xi\vec x$：
\[\chi_T\propto\xi^{2-\eta}\int d^dx\,\frac{f(|\vec x|)}{|\vec x|^{d-2+\eta}}\propto|T-T_c|^{-\nu(2-\eta)}\qquad\Rightarrow\qquad\gamma=\nu(2-\eta).\tag{9.27}\]
〔注〕書上 (9.27) 的被積函數只寫 $f(|\vec x|)$，少了 $|\vec x|^{-(d-2+\eta)}$，不影響結論。一維核對：$\chi/\xi=2.0123,\ 2.0002,\ 2.0000$（$K=1,2,3$），即 $\chi\propto\xi^{2-\eta}$、$\eta=1$
（一維 $T_c=0$，臨界點上 $G\equiv1$ 不衰減，$d-2+\eta=0$）。

\textbf{四條標度律}（(9.16)--(9.19)）：
\[\alpha+2\beta+\gamma=2,\qquad\gamma=\nu(2-\eta),\qquad\alpha=2-d\nu,\qquad\delta=\frac{d+2-\eta}{d-2+\eta}.\tag{9.16--9.19}\]
平均場、二維 Ising、二維自迴避三組精確指數每一條都成立（圖 \ref{nine-fig-exp} 右，用分數驗算）；平均場的後兩條（hyperscaling）只在 $d=4$ 成立（$d=3$ 時 $2-d\nu=\frac12\ne0$）。
六個指數只有兩個獨立：例如知道 $\nu,\eta$，就有 $\gamma=(2-\eta)\nu$、$\alpha=2-d\nu$、$\beta=\nu(d-2+\eta)/2$ 與 $\delta$。

\textbf{Widom 的齊次形式}（(9.20)--(9.23)，練習 9.1--9.2）：靠近臨界點，自由能只透過一個組合依賴 $h$，
\[G(T,h)=|r|^{2-\alpha}f\Big(\frac h{|r|^\Delta},\frac r{|r|}\Big),\qquad r=\frac{T-T_c}{T_c}.\tag{9.21, 9.23}\]
對 $h$ 微分一次、兩次得 $\beta=2-\alpha-\Delta$、$\gamma=\Delta-\beta$（合起來就是 (9.16)）；$r=0$ 時得 $\delta=\Delta/\beta$，即 $\beta(\delta-1)=\gamma$（Widom）。二維 Ising $\alpha=0$、$\Delta=\frac{15}8$：
$\beta=\frac18$、$\gamma=\frac74$、$\delta=15$；用一個具體的 $f$ 數值微分，斜率量到 0.1250、1.7500、0.0667（$=\frac1{15}$），而 $\frac18\times14=\frac74$。

\begin{figure}[htbp]
\centering
\begin{minipage}[c]{0.5\linewidth}\centering\footnotesize\setlength{\tabcolsep}{2pt}\NineTabNineOne\\[8pt]\NineFigChain\end{minipage}\hfill
\begin{minipage}[c]{0.48\linewidth}\centering\footnotesize\setlength{\tabcolsep}{2.2pt}\NineTabLaws\end{minipage}
\caption{臨界指數。左上：書上的 Table 9.1（紅字是筆誤、綠字是更正；$\gamma'$ 無定義是因為 $\chi_T=\infty$）。左下：一維 Ising 鏈的 $G(r)=t^{|r|}$，正好落在 $e^{-|r|/\xi}$ 上。
右：三組精確指數與書上的三維欄，逐格驗算標度律；書上三維欄兩格不成立，都來自 $\alpha$}\label{nine-fig-exp}
\end{figure}

\textbf{揭曉}：六個指數只有兩個獨立；知道 $\nu$ 和 $\eta$，其餘都由標度律決定。

% =====================================================================
\subsection{Ginzburg--Landau--Wilson：讓自旋可以伸縮}
\leavevmode\bk{pp.~215--217}\ \anim{W904}\par\nopagebreak

\textbf{問}：如果自旋不是固定長度的箭頭，而是可以伸縮的「場」，磁鐵的自由能長什麼樣子？

\textbf{有效 Hamiltonian}（(9.28)(9.29)）：
\[\mathcal H_{GL}=\sum_i\Big\{\frac r2|\vec S_i|^2+\frac12|\nabla\vec S_i|^2+\frac u4|\vec S_i|^4-\vec h\cdot\vec S_i\Big\},\qquad Z=\prod_i\int d^nS_i\,e^{-\mathcal H_{GL}}.\tag{9.28, 9.29}\]
三項各有用意。$|\nabla\vec S|^2$（相鄰格點之差的平方）取代 $-J\vec S_i\cdot\vec S_j$：因為 $-\vec S_i\cdot\vec S_j=\frac12|\vec S_i-\vec S_j|^2-\frac12(|\vec S_i|^2+|\vec S_j|^2)$，鄰居想要相等。
$\frac r2|\vec S|^2+\frac u4|\vec S|^4$ 是長度的「軟約束」，$r\propto T-T_c$ 是約化溫度。指數裡沒有 $k_BT$：溫度已經吸收進係數，所以叫「有效」Hamiltonian。

\textbf{Fourier 分解與平均場}（(9.30)--(9.34)，練習 9.3）：$\vec S(\vec q)=\frac1N\sum_i\vec S_ie^{i\vec q\cdot\vec r_i}$ 讓二次項對各個 $\vec q$ 分開成 $(r+|\vec q|^2)|\vec S(\vec q)|^2$（格子上 $|\vec q|^2$ 是 $4\sin^2\frac q2$），
四次項把 $\vec q_1+\vec q_2+\vec q_3+\vec q_4=0$ 的四個模態連起來；8 點環（種子 904）的離散版，三個恆等式的誤差都 $<10^{-12}$。〔注〕(9.32) 的 $V$ 與 $(2\pi)^d$ 因子和 (9.30)(9.31) 的正規化對不上，不影響 (9.33)。
\textbf{只留 $\vec q=0$}（所有 $\vec S_i$ 都等於 $\vec m$）就是平均場：
\[\mathcal H_{MF}=V\Big[\frac r2|\vec m|^2+\frac u4|\vec m|^4-\vec h\cdot\vec m\Big],\qquad Z_{MF}=\int d^nm\,e^{-\mathcal H_{MF}(\vec m)}.\tag{9.33, 9.34}\]

\textbf{算到底：雙井}（$n=1$、$u=1$，圖 \ref{nine-fig-glw} 左）：$F(m)=\frac r2m^2+\frac u4m^4-hm$，$F'(m)=rm+um^3-h=0$。$h=0$ 時 $r>0$ 只有 $m=0$ 一個谷；$r<0$ 有兩個谷 $m^*=\pm\sqrt{-r/u}$，
$F_{\min}=\frac r2\cdot\frac{-r}u+\frac u4\cdot\frac{r^2}{u^2}=-\frac{r^2}{4u}$（$r=-1$：$m^*=\pm1$、$F_{\min}=-0.25$）。小球滾進其中一個谷＝對稱破缺；$h=0.1$ 把雙井歪向右，谷底移到 $m^*=1.0467$。谷底 $\propto(-r)^{1/2}$，所以平均場 $\beta=\frac12$。

\textbf{練習 9.4}（書上 (9.35) $F=\frac r2m^2+\frac{m^4}{12}-hm$，即 $u=\frac13$）：$r=-1$ 時 $m^*=\sqrt3=1.7321$、$F_{\min}=-0.75$。令 $m=|r|^{1/2}\mathcal M$、$h=x|r|^{3/2}$，
方程變成與 $|r|$ 無關的 $\frac{\mathcal M^3}3-\mathcal M=x$（$x=1$：$\mathcal M=2.1038$，$r=-0.01$ 與 $-0.0001$ 相同），所以 $\beta=\frac12$、$\Delta=\frac32$；$F_{\min}\propto|r|^2$ 給 $\alpha=0$；
另有 $\delta=3$（$r=0$：$m=(3h)^{1/3}$）、$\gamma=1$（數值斜率見圖 \ref{nine-fig-glw} 右下）。〔注〕(9.35) 的四次項是 $\frac{m^4}{12}$，不是 (9.33) 的 $\frac u4m^4$ 配 $u=1$，指數不受影響；(b) 小題的 $r^{2-\alpha}$ 應為 $|r|^{2-\alpha}$。

\textbf{有限體積}（(9.34)，$r=-1$、$u=1$；圖 \ref{nine-fig-glw} 右上）：$h=0.01$ 時 $\langle m\rangle=0.0869,\ 0.7576,\ 1.0042$（$V=10,100,1000$），趨近谷底 $m^*=1.0050$；
$h=0$ 時不論 $V$ 多大 $\langle m\rangle=0$，但 $\langle|m|\rangle=0.8888,\ 0.9922,\ 0.9992$；$-\ln Z/V=-0.2719,\ -0.2397,\ -0.2478\to F_{\min}=-0.25$。
先 $V\to\infty$、再 $h\to0$ 才得到 $m^*\ne0$——和 9.2 節的翻轉等待是同一件事。

\begin{figure}[htbp]
\centering
\begin{minipage}[c]{0.5\linewidth}\centering\NineFigWell\end{minipage}\hfill
\begin{minipage}[c]{0.48\linewidth}\centering\footnotesize\setlength{\tabcolsep}{3pt}\NineTabZmf\\[8pt]\NineTabExNineFour\end{minipage}
\caption{GLW 的平均場。左：雙井 $F=\frac r2m^2+\frac14m^4-hm$，$r=1,0,-1$（藍，由淺到深）；$r=-1$ 的兩個谷 $m^*=\pm1$（綠點），外場 $h=0.1$ 把它歪向右（橙虛線，球在新的谷底）。
右上：(9.34) 的有限體積平均。右下：練習 9.4 的數值斜率}\label{nine-fig-glw}
\end{figure}

\textbf{揭曉}：GLW 把自旋變成可伸縮的場；只留 $\vec q=0$ 就是平均場：雙井給出 $\beta=\frac12$、$\delta=3$、$\alpha=0$。

% =====================================================================
\subsection{從磁鐵讀出漫步：\texorpdfstring{$\langle R^2\rangle\propto N^{2\nu}$}{⟨R²⟩ 正比於 N 的 2ν 次方}}
\leavevmode\bk{pp.~218--219}\ \anim{W905}\par\nopagebreak

\textbf{問}：只知道磁鐵的 $\xi\propto(z_c-z)^{-\nu}$、$\chi\propto(z_c-z)^{-\gamma}$，能不能推出漫步的 $\langle R^2\rangle\propto N^{2\nu}$？

\textbf{先借用 9.6 節的結論}（(9.45)）：$n\to0$ 的磁鐵，關聯函數就是自迴避漫步的生成函數 $G(z;\vec R)=\sum_NC(N;\vec R)z^N$
（$C$＝從原點走 $N$ 步到 $\vec R$ 的自迴避條數，$z$ 是每步的逸度，$z\to z_c$ 就是 $T\to T_c$）。於是
\[\xi^2=\frac{\sum_{\vec R}R^2G}{\sum_{\vec R}G}=\frac{\sum_Nc_N\langle R^2\rangle_Nz^N}{\sum_Nc_Nz^N},\qquad\chi=\sum_Nc_Nz^N .\tag{9.36, 9.37}\]
分子把每一條漫步的 $R^2$ 加起來，分母數條數。\textbf{算到底}（二維，第 7 章的精確計數）：
\begin{center}
\footnotesize\setlength{\tabcolsep}{3.5pt}\NineTabSaw\hfill\NineTabCoef
\end{center}
$N=2$：直走 4 條（$R^2=4$）、轉彎 8 條（$R^2=2$），$16+16=32$。$z=0.1$（$N\le14$ 的部分和）：分子 0.9964、$\chi=1.5699$，$\xi^2=0.6347$；$z=0.15,\ 0.2$：$\xi^2=1.2799,\ 2.4593$——$z$ 越接近 $z_c=0.3791$，$\xi$ 越大。
〔注〕書上寫 $\Gamma_N(z/z_0)^N$，$z_0=(2dJ)^{-1}$ 是平均場的臨界逸度，所以奇點在 $z_c=z_0/\mu$；用每步的逸度就是 $z_c=1/\mu$。

\textbf{取係數}（第 2 章補充 3 的 (2:S3-10)；〔注〕書上說「Chapter 1」）：
\[[z^N](z_c-z)^{-s}=z_c^{-N-s}\frac{\Gamma(N+s)}{\Gamma(s)N!}\approx\frac{N^{s-1}}{\Gamma(s)}z_c^{-(N+s)} .\]
精確 $\div$ 漸近見上右表（$s=\gamma$ 與 $s=\gamma+2\nu$），$N$ 越大越接近 1；$s$ 越大收斂越慢。\textbf{兩個係數一除}（(9.38)--(9.41)）：$\xi^2\chi=K_1K_2(z_c-z)^{-(\gamma+2\nu)}$ 的係數是 $c_N\langle R^2\rangle_N$，$\chi=K_2(z_c-z)^{-\gamma}$ 的係數是 $c_N$，
\[\langle R^2\rangle_N=\frac{K_1K_2N^{\gamma+2\nu-1}z_c^{-(N+\gamma+2\nu)}/\Gamma(\gamma+2\nu)}{K_2N^{\gamma-1}z_c^{-(N+\gamma)}/\Gamma(\gamma)}=\frac{K_1}{z_c^{2\nu}}\frac{\Gamma(\gamma)}{\Gamma(\gamma+2\nu)}N^{2\nu}.\tag{9.41}\]
兩個精確係數之比 $\div$ 這個式子：1.2476、1.0240、1.0024 $\to1$（上右表最後一列）。
\begin{errbox}[p.~218--219，(9.37)--(9.41)]
(1) (9.37) 之後說關聯長度 $\propto K_1/(z_c-z)^\nu$；要得到 (9.38) 的 $K_1K_2$，應是 $\xi^2\approx K_1/(z_c-z)^{2\nu}$（否則後面的 $K_1$ 都要換成 $K_1^2$）。
(2) (9.38)$\to$(9.39) 的係數印成 $z_c^{-(N+s)}N^{\delta-1}/\Gamma(s)$，應為 $N^{s-1}$；照書上（$\delta=\frac{91}5$）$N=10$ 的比值是 $1.4\times10^{-17}$，完全不對。
(3) (9.40) 之後「must be equal to $V\Gamma_N/z_0$」應為 $V\Gamma_N/z_0^N$。
(4) (9.41) 印成 $K_1/z_c^{\gamma}$，應為 $K_1/z_c^{2\nu}$：$z_c^{-(N+\gamma+2\nu)}\div z_c^{-(N+\gamma)}=z_c^{-2\nu}$；二維 $\mu^{3/2}=4.2850$，書上給 $\mu^{43/32}=3.6823$。
\end{errbox}
\textbf{回到第 7 章}（$N=10,12,14$）：$2\nu_{\rm eff}=1.4442,\ 1.4505,\ 1.4556\to\frac32$；$\langle R^2\rangle_N/N^{3/2}=0.8299,\ 0.8224,\ 0.8168$，慢慢趨於常數：前因子與 $N$ 無關，$\langle R^2\rangle\propto N^{2\nu}$——de Gennes。

\textbf{揭曉}：磁鐵的 $\xi$ 和 $\chi$ 只要是冪次，取 $z^N$ 的係數一除，就得到漫步的 $\langle R^2\rangle\propto N^{2\nu}$。

% =====================================================================
\subsection{\texorpdfstring{$n\to0$}{n→0} 的魔術：O(\texorpdfstring{$n$}{n}) 模型就是自迴避漫步}
\leavevmode\bk{pp.~219--231}\ \anim{W906}\par\nopagebreak

\textbf{問}：「$n=0$ 個分量的磁鐵」聽起來荒謬——它為什麼會數自迴避漫步？目標是 (9.45)：
\[\lim_{n\to0}\frac1n\sum_{\alpha=1}^n\langle S_{i\alpha}S_{j\alpha}\rangle=G(z;\vec R_i,\vec R_j)\equiv\sum_NC(N;\vec R_i,\vec R_j)z^N .\tag{9.45}\]

\textbf{先在小圖上看見它}（圖 \ref{nine-fig-n0} 左）：$2\times3$ 梯子（6 點 7 鍵），權重 $\prod_{\text{鍵}}(1+K\vec S_i\cdot\vec S_j)$，自旋長度 $|\vec S|^2=n$（每個分量 $\langle S_\alpha^2\rangle=1$）。
乘開：每條鍵「不選」（1）或「選」（$K\vec S_i\cdot\vec S_j$）。對自旋方向平均時，一個格點上的 $\vec S$ 出現奇數次就平均成 0，所以只有「每點用 0 或 2 次」的鍵集合活下來（這張圖每點最多 3 條鍵）；
一個閉圈 $(\vec S_1\cdot\vec S_2)(\vec S_2\cdot\vec S_3)\cdots(\vec S_k\cdot\vec S_1)$ 平均後是 $\sum_\alpha1=n$。分母 $Z$ 是閉圈的和；分子 $\langle S_{0\alpha}S_{3\alpha}\rangle$ 還要一條從 0 到 3 的路（加上不相交的閉圈）：
\[G_{03}=\frac{K+K^3+K^5+nK^5}{1+2nK^4+nK^6},\qquad G_{05}=\frac{3K^3+K^5}{1+2nK^4+nK^6}.\]
$K=0.5$：$G_{03}=0.6027$（$n=1$，64 個 Ising 組態暴力加總）、0.5610（$n=2$）。$n\to0$ 時\textbf{每個閉圈都帶一個 $n$ 而消失}，分母變 1，剩下 $K+K^3+K^5=0.65625$
——正是三條從 0 到 3 的自迴避路徑（長 1、3、5）；對角的 $G_{05}$：0.3562、0.3171、0.40625$\,=3K^3+K^5$（長 3、3、3、5）。

\textbf{書上的路：六步走完 (9.42)--(9.77)}（只看主軸，(9.47)--(9.65) 的細節不逐行核對）：
\begin{enumerate}[label=\textbf{\arabic*.},itemsep=0.1em,leftmargin=1.6em]
\item \textbf{加源}（(9.42)--(9.44)）：在 GLW 的指數裡加 $\sum_i\vec h_i\cdot\vec S_i$，$\partial^2\ln Z/\partial h_{i\alpha}\partial h_{j\alpha}$ 就是關聯函數。〔注〕(9.44) 左邊書上寫 $\partial^2Z$，應為 $\partial^2\ln Z$。
\item \textbf{四次項變二次}（(9.46)）：每個格點引進一個輔助場 $\phi_i$，
\[e^{-\frac u4|\vec S_i|^4}=\frac1{\sqrt{\pi u}}\int_{-\infty}^\infty d\phi_i\,e^{-\phi_i^2/u-i\phi_i|\vec S_i|^2}\tag{9.46}\]
（一維高斯積分；$u=1$、$|\vec S|^2=1,2$ 的數值積分得 0.7788、0.3679，就是 $e^{-1/4}$、$e^{-1}$）。
\item \textbf{每個分量一樣}（(9.47)--(9.62)）：固定 $\phi$，$\vec S$ 的積分是高斯積分，$n$ 個分量各給一個相同的因子 $(\det\overleftrightarrow V)^{-1/2}$，$\overleftrightarrow V=(\frac r2-\frac12\nabla^2)\overleftrightarrow I+i\overleftrightarrow\phi$；
所以 $Z_n\propto(\det\overleftrightarrow V)^{-n/2}\to1$（$n\to0$），而 $\langle S_{i\alpha}S_{j\alpha}\rangle\propto(\overleftrightarrow V^{-1})_{ij}$（本節末的補充），再對 $\phi$ 以權重 $e^{-\sum\phi_i^2/u}$ 平均。$n\to0$ 的作用就是讓分母消失。
\item \textbf{逆矩陣寫成積分}（(9.63)--(9.65)，練習 9.5）：$\overleftrightarrow V^{-1}=\int_0^\infty e^{-t\overleftrightarrow V}dt$，所以
\[G(\vec R_i,\vec R_j)=\prod_l\int\frac{d\phi_l}{\sqrt{\pi u}}e^{-\sum\phi_l^2/u}\int_0^\infty dt\,e^{-rt/2}\langle\vec R_i|e^{(\frac12\nabla^2-i\phi)t}|\vec R_j\rangle .\tag{9.65}\]
〔注〕書上的 (9.65) 沒寫出 $\phi$ 的高斯權重 $e^{-\sum\phi_i^2/u}$（(9.62) 第二行有）。
\item \textbf{切片成路徑積分}（(9.66)--(9.71)，練習 9.6）：$t$ 切成 $M$ 段 $\epsilon$，$e^{(\frac12\nabla^2-i\phi)\epsilon}\approx e^{-i\phi\epsilon}e^{\frac12\nabla^2\epsilon}$，段與段之間插入 $\sum_{\vec R}|\vec R\rangle\langle\vec R|$：
$e^{-i\phi\epsilon}$ 是對角的（記下路徑經過處的 $\phi$），$e^{\frac12\nabla^2\epsilon}$ 是擴散核（第 2 章的高斯）。誤差（6 點環、種子 906，圖 \ref{nine-fig-n0} 右下）：一階 $\propto\epsilon^2$，乘上交換子項後 $\propto\epsilon^3$；
$t=1$ 的總誤差 $\propto1/M$，$M\to\infty$ 才精確。
\item \textbf{對 $\phi$ 積分}（(9.72)--(9.77)）：一個格點被走過 $k$ 次，$\phi$ 的積分給 $\frac1{\sqrt{\pi u}}\int d\phi\,e^{-\phi^2/u-ik\phi}=e^{-uk^2/4}$：$k=1$ 給 $e^{-u/4}$、$k=2$ 給 $e^{-u}$，
比「兩次分開經過」多了 $e^{-u}/(e^{-u/4})^2=e^{-u/2}$（$u=1$：0.6065）——\textbf{每相交一次乘 $e^{-u/2}$}（$k^2=k+2\binom k2$）。於是 $G=\int_0^\infty dt\,e^{-rt/2}C(t;\vec R_i,\vec R_j)$
（$C$＝相交受罰的漫步之和）；換成步數，$G=\sum_NC(N)e^{-rN/2}$、$e^{-r/2}=z/z_c$（(9.76)(9.77)）：就是 (9.45)。
\end{enumerate}
\begin{errbox}[p.~225，(9.70)]
書上寫 $\langle\vec R_M|e^{\nabla^2\epsilon}|\vec R_{M'}\rangle=(2\pi)^{-d/2}\exp\big(-\frac d2|\vec R_M-\vec R_{M'}|^2/\epsilon\big)$；應為
$\langle\vec R_M|e^{\frac12\nabla^2\epsilon}|\vec R_{M'}\rangle=(2\pi\epsilon)^{-d/2}\exp\big(-|\vec R_M-\vec R_{M'}|^2/2\epsilon\big)$。照書上，核的總積分是 $(\epsilon/d)^{d/2}$（$\epsilon=0.1$：一維 0.3162、三維 0.00609），不是 1。
\end{errbox}
\begin{errbox}[p.~226，(9.72)(9.75)]
(1) (9.72) 的積分下限是 0，應為 $\int_{-\infty}^\infty$：從 0 積得 $0.3894-0.2395i$，不是 $e^{-1/4}=0.7788$。
(2) (9.75) 指數裡兩項之間少了負號，應為 $-\int\frac d2|\frac{d\vec R}{ds}|^2ds-\frac u4\iint\delta(\vec R(s')-\vec R(s))\,ds\,ds'$；照書上每相交一次乘 $e^{+u/2}$，越相交越多（圖 \ref{nine-fig-n0} 左下最後一欄）。
\end{errbox}

\textbf{弱自迴避}（圖 \ref{nine-fig-n0} 左下）：$N=8$ 的全部 $4^8$ 條漫步，每相交一次乘 $e^{-u/2}$：$u=0,1,4,\infty$ 得 65536、23695.11、7988.15、5916——$u\to\infty$ 只剩自迴避的 $c_8=5916$；照書上的正號，$u=1$ 得 $6.2\times10^5$。
〔注〕所以 $C(t;\cdots)$ 是「弱自迴避」（Edwards 模型）：$u$ 有限時相交只是受罰，$u\to\infty$ 才嚴格；兩者在臨界點附近屬於同一個普適類，指數相同。
p.~227 寫「$r/2=1-T_c/T$、$e^{-r/2}\approx T_c/T=J/T=z/z_c$」：和 (9.21) 的 $r=(T-T_c)/T_c$ 差 2 倍（只差常數倍，指數不受影響）；中間的「$=J/T$」只在 $z_c=1$ 時成立。

\textbf{補充（書上 9.6 節）：高斯積分}（(9:S-1)--(9:S-22)）。轉到 $\hat T$ 的本徵座標，$N$ 維積分變成 $N$ 個一維高斯積分的乘積：
\[\prod_i\int dx_i\,e^{-\frac12\sum x_iT_{ij}x_j+\sum h_ix_i}=\frac{(2\pi)^{N/2}}{\sqrt{\det\hat T}}\exp\Big(\frac12\sum_{i,j}h_i(\hat T^{-1})_{ij}h_j\Big),\qquad\langle x_rx_s\rangle=(\hat T^{-1})_{rs}.\tag{9:S-11, 9:S-18}\]
例：$\hat T=\begin{pmatrix}2&-1\\-1&2\end{pmatrix}$、$\vec h=(1,0)$：$\det\hat T=3$、$\hat T^{-1}=\frac13\begin{pmatrix}2&1\\1&2\end{pmatrix}$，積分 $=\frac{2\pi}{\sqrt3}e^{1/3}=5.0627$（二重數值積分相同）；$\langle x_1x_2\rangle=\frac13$、$\langle x_1^2\rangle=\frac23$。
四次動差用 \textbf{Wick 定理} (9:S-20)：對所有配對的乘積求和，$\langle x_1^2x_2^2\rangle=\langle x_1^2\rangle\langle x_2^2\rangle+2\langle x_1x_2\rangle^2=\frac49+\frac29=\frac23$、$\langle x_1^4\rangle=3\cdot\frac49=\frac43$；
$2m$ 個變數的配對數 $\frac{(2m)!}{2^mm!}=1,3,15,105$（第 7 章的弧線圖）。第 3 步用的就是 (9:S-11)：每個分量一個 $(\det)^{-1/2}$。〔注〕(9:S-16) 左邊的 $\langle x_rx_k\rangle$ 應為 $\langle x_rx_s\rangle$。
\begin{errbox}[p.~230，(9:S-19)]
第二行寫 $\frac{\partial^3}{\partial x_r\partial x_s\partial x_t}$，應為 $\frac{\partial^3}{\partial h_r\partial h_s\partial h_t}$；第三行括號裡的 $\langle x_rx_s\rangle$ 應為 $\langle x_tx_u\rangle$（對 $h_u$、$h_t$ 微分後留下的是 $\langle x_tx_u\rangle$）。
照書上再微分，$\hat T=\begin{pmatrix}2&-1\\-1&3\end{pmatrix}$ 的 $\langle x_1^2x_2^2\rangle$ 得 0.44，正確是 0.32。
\end{errbox}

\begin{figure}[htbp]
\centering
\begin{minipage}[c]{0.5\linewidth}\centering\NineFigLadder\\[6pt]\footnotesize\setlength{\tabcolsep}{2.5pt}\NineTabWeak\end{minipage}\hfill
\begin{minipage}[c]{0.48\linewidth}\centering\footnotesize\setlength{\tabcolsep}{2.5pt}\NineTabLadder\\[8pt]\setlength{\tabcolsep}{3pt}\NineTabTrotter\end{minipage}
\caption{$n\to0$。左上：$2\times3$ 梯子乘開後存活的項（藍點是端點，橙色是路徑，紅色是帶一個 $n$ 的閉圈）；$n\to0$ 時紅色的全部消失，分母變成 1。
左下：$N=8$ 的弱自迴避計數。右上：關聯函數（$n=1$ Ising、$n=2$ xy）。右下：Trotter 分解的誤差（練習 9.6）}\label{nine-fig-n0}
\end{figure}

\textbf{揭曉}：$n\to0$ 時每個閉圈都帶一個 $n$ 而消失，剩下的就是一條從 $i$ 到 $j$、不碰自己的路——O($n$) 的關聯函數就是自迴避漫步的生成函數。

% =====================================================================
\subsection{練習 9.1--9.6 解答}
{\small
\begin{enumerate}[label=\textbf{9.\arabic*},itemsep=0.05em]
\item Widom 形式 $\Rightarrow$ (9.16)。\quad 答：$M=-\partial G/\partial h=-|r|^{2-\alpha-\Delta}f_1(0,\pm1)$（$f_1$＝對第一個變數的導數）$\Rightarrow\beta=2-\alpha-\Delta$；
$\chi=\partial M/\partial h\propto|r|^{2-\alpha-2\Delta}\Rightarrow\gamma=2\Delta-2+\alpha=\Delta-\beta$。消去 $\Delta$ 得 $\alpha+2\beta+\gamma=2$。二維 Ising（$\alpha=0$、$\Delta=\frac{15}8$）：$\beta=\frac18$、$\gamma=\frac74$；
取 $f(x)=-(1+x^2)^{q/2}-x(1+x^2)^{(q-1)/2}$（$q=\frac{2-\alpha}\Delta$）數值微分，斜率 0.1250、1.7500。
\item Widom 定律。\quad 答：（題目的「$t=0$」即 $r=0$。）$r\to0$ 時 $G$ 要有限且與 $r$ 無關 $\Rightarrow$ 大 $x$ 時 $f\sim x^{(2-\alpha)/\Delta}$、$G\propto h^{(2-\alpha)/\Delta}$ $\Rightarrow$ $M\propto h^{(2-\alpha-\Delta)/\Delta}=h^{\beta/\Delta}$，
$\delta=\Delta/\beta$，$\beta(\delta-1)=\Delta-\beta=\gamma$。平均場（$\Delta=\frac32$）3、二維 Ising（$\frac{15}8$）15、二維自迴避（$\frac{91}{64}$）$\frac{91}5$；上例 $1/\delta$ 量到 0.0667。
\item GLW 的 Fourier 形式。\quad 答：代入 (9.31)，用 $\sum_ie^{i(\vec q+\vec q\,')\cdot\vec r_i}=(2\pi)^d\delta(\vec q+\vec q\,')$（格距 1、$V=N$）：二次項 $\frac{V^2}{2(2\pi)^d}\int(r+|\vec q|^2)|\vec S(\vec q)|^2d^dq$（格子上 $|\vec q|^2\to\sum_k4\sin^2\frac{q_k}2$），
四次項 $\frac{uV^4}{4(2\pi)^{3d}}\int(\vec S_1\cdot\vec S_2)(\vec S_3\cdot\vec S_4)\delta(\vec q_1+\cdots+\vec q_4)\prod_kd^dq_k$，外場 $-V\vec h\cdot\vec S(0)$；書上 (9.32) 只差常數因子。離散版（8 點環）誤差 $<10^{-12}$。
\item Landau 自由能。\quad 答：$rm+\frac{m^3}3-h=0$；$h=0$、$r<0$：$m^*=\sqrt{-3r}$、$F_{\min}=-\frac34r^2$。令 $m=|r|^{1/2}\mathcal M$、$h=x|r|^{3/2}$ 得 $\frac{\mathcal M^3}3\pm\mathcal M=x$（與 $|r|$ 無關；$x=1$、$r<0$：2.1038）
$\Rightarrow\beta=\frac12$、$\Delta=\frac32$；$F_{\min}=|r|^2\mathcal K\Rightarrow\alpha=0$；$\delta=3$；$\chi=1/r$（$r>0$）、$1/2|r|$（$r<0$，$F''(m^*)=-2r$）$\Rightarrow\gamma=\gamma'=1$，振幅比 2。
\item 逆矩陣的積分表示。\quad 答：$\overleftrightarrow V\int_0^\infty e^{-t\overleftrightarrow V}dt=\int_0^\infty\big(-\frac d{dt}e^{-t\overleftrightarrow V}\big)dt=\overleftrightarrow I-\lim_{t\to\infty}e^{-t\overleftrightarrow V}=\overleftrightarrow I$，因為
$\mathrm{Re}\,\langle x,\overleftrightarrow Vx\rangle=\langle x,(\frac r2-\frac12\nabla^2)x\rangle\ge\frac r2|x|^2$（$i\phi$ 只貢獻虛部），$|e^{-t\overleftrightarrow V}x|\le e^{-rt/2}|x|\to0$。$\overleftrightarrow V$ 與 $e^{-t\overleftrightarrow V}$ 對易，左逆也是右逆；數值（6 點環）誤差 $<10^{-10}$。
\item Trotter 公式的下一項。\quad 答：$e^{\epsilon(X+Y)}=e^{\epsilon X}e^{\epsilon Y}e^{-\frac{\epsilon^2}2[X,Y]+O(\epsilon^3)}$（兩邊展開到 $\epsilon^2$ 比較）；$X=-i\phi$、$Y=\frac12\nabla^2$：$-\frac{\epsilon^2}2[X,Y]=\frac{i\epsilon^2}4[\phi,\nabla^2]$。
一階誤差相鄰比 3.85、3.92 $\to4$，加上這一項後 7.63、7.81 $\to8$；$t=1$、$M=10,20,40$ 段的總誤差 0.02366、0.0117、0.005821（$\propto1/M$）。
\end{enumerate}}

% =====================================================================
\subsection{第 9 章的回顧與筆誤一覽}
\begin{center}
\small\setlength{\tabcolsep}{4pt}
\begin{tabular}{@{}llll@{}}
\toprule
主題 & 結果 & 書上的式子 & 動畫\\
\midrule
臨界指數 & $p=2\nu=\frac6{d+2}$；$M_s\propto(T_c-T)^\beta$：平均場 $\frac12$、二維 Ising $\frac18$ & (9.1)--(9.3) & W901\\
對稱破缺 & $\langle M\rangle=0$ 但 $\langle|m|\rangle=0.9337$；翻轉等待 133、608、2151 & (9.4)--(9.15) & W902\\
標度律 & 一維 $\chi=e^{2K}$、$\eta=1$；四條標度律；三維 $\alpha=0.12\to0.19$ & (9.16)--(9.27) & W903\\
GLW & $m^*=\sqrt{-r/u}$；$\beta=\frac12$、$\Delta=\frac32$、$\alpha=0$、$\delta=3$ & (9.28)--(9.35) & W904\\
磁鐵 $\to$ 漫步 & $\langle R^2\rangle_N=\frac{K_1}{z_c^{2\nu}}\frac{\Gamma(\gamma)}{\Gamma(\gamma+2\nu)}N^{2\nu}$ & (9.36)--(9.41) & W905\\
$n\to0$ & 每個閉圈 $\times n$、每相交 $e^{-u/2}$；$G=\sum_NC(N)(z/z_c)^N$ & (9.42)--(9.77) & W906\\
\bottomrule
\end{tabular}
\end{center}
\textbf{筆誤一覽}（給 2004 年版；〔注〕見各節）：下表在文中都用紅框標出，每一條都用程式核對過。
\begin{center}
\footnotesize\renewcommand{\arraystretch}{0.9}
\NineTabErrata
\end{center}
